% Generated by replication/build_presentation_assets.py
\begin{frame}{Paper Table 1}
\begin{minipage}{\textwidth}\small\textbf{Employment shares and mean real hourly income, 2008 vs. 2025}\end{minipage}\par\vspace{0.15cm}
\centering\begingroup\fontsize{8.5}{10}\selectfont
\setlength{\tabcolsep}{4pt}\renewcommand{\arraystretch}{1.06}
\begin{adjustbox}{max width=\textwidth,max totalheight=0.64\textheight}
% Generated from Paper/sections/descriptive_core_size_2008_2025_table.tex
\begin{tabular}{lrrrrrrr}
\toprule
    & \multicolumn{2}{c}{Employment share} & \multicolumn{2}{c}{Mean hourly income} & \multicolumn{2}{c}{Wage ratio to solo} & \multicolumn{1}{c}{Wage growth} \\
    \cmidrule(lr){2-3} \cmidrule(lr){4-5} \cmidrule(lr){6-7} \cmidrule(lr){8-8}
    Firm size & 2008 & 2025 & 2008 & 2025 & 2008 & 2025 & Annualized \\
    \midrule
Solo & 36.4\% & 34.6\% & 5,179 & 6,394 & 1.0 & 1.0 & 1.2\% \\
2-3 & 19.1\% & 13.7\% & 5,871 & 7,097 & 1.1 & 1.1 & 1.1\% \\
4-5 & 7.4\% & 7.0\% & 6,790 & 8,096 & 1.3 & 1.3 & 1.0\% \\
6-10 & 6.0\% & 5.8\% & 7,481 & 9,426 & 1.4 & 1.5 & 1.4\% \\
11-19 & 3.5\% & 3.9\% & 8,825 & 11,446 & 1.7 & 1.8 & 1.5\% \\
20-30 & 3.3\% & 3.3\% & 9,175 & 10,990 & 1.8 & 1.7 & 1.1\% \\
31-50 & 2.6\% & 3.0\% & 10,862 & 12,408 & 2.1 & 1.9 & 0.8\% \\
51-100 & 3.1\% & 3.5\% & 9,918 & 13,039 & 1.9 & 2.0 & 1.6\% \\
101+ & 18.5\% & 25.2\% & 12,757 & 16,549 & 2.5 & 2.6 & 1.5\% \\
\addlinespace
    All workers & 100.0\% & 100.0\% & 7,525 & 10,103 & 1.5 & 1.6 & 1.7\%\\
\bottomrule\end{tabular}

\end{adjustbox}\endgroup\par\vspace{0.10cm}
{\fontsize{6.5}{7.5}\selectfont\raggedright Mean hourly income is expressed in constant 2025 pesos. Annualized income growth is compounded over 2008--2025.\par}
\note{Mean hourly income is expressed in constant 2025 pesos. Annualized income growth is compounded over 2008--2025.}
\end{frame}

\begin{frame}{Paper Table 2 (1/3)}
\begin{minipage}{\textwidth}\small\textbf{Mean real hourly labor income gaps by worker group, 2025}\end{minipage}\par\vspace{0.15cm}
\centering\begingroup\fontsize{8.5}{10}\selectfont
\setlength{\tabcolsep}{4pt}\renewcommand{\arraystretch}{1.06}
\begin{adjustbox}{max width=\textwidth,max totalheight=0.64\textheight}
% Generated from Paper/sections/descriptive_wage_gaps_2025_table.tex
\begin{tabular}{lp{0.34\textwidth}rrrr}
\toprule
    Dimension & Group & Employment share & Mean income & Gap & Ratio \\
    \midrule
Sex & Men & 58.8\% & 10,236 & -0.6\% & 1.00 \\
& Women & 41.2\% & 10,379 & 0.8\% & 1.01 \\
\midrule
    Formality & Formal & 44.9\% & 15,241 & 48.0\% & 2.43 \\
& Informal & 55.1\% & 6,271 & -39.1\% & 1.00 \\
\midrule
    Education & Preschool or less & 2.9\% & 4,333 & -57.9\% & 1.00 \\
& Primary & 18.1\% & 5,693 & -44.7\% & 1.31 \\
& Secondary & 45.5\% & 7,085 & -31.2\% & 1.63 \\
& Higher education & 33.5\% & 17,667 & 71.6\% & 4.08\\
\bottomrule\end{tabular}

\end{adjustbox}\endgroup\par\vspace{0.10cm}
{\fontsize{6.5}{7.5}\selectfont\raggedright Employment shares are computed within each dimension. Mean income is real hourly labor income in constant 2025 pesos. The gap is the percent difference in each group's mean relative to the overall 2025 mean. The ratio divides each group's mean by the lowest mean within the same dimension. The sample excludes employed pension recipients. Sector rows omit extraterritorial organizations.\par}
\note{Employment shares are computed within each dimension. Mean income is real hourly labor income in constant 2025 pesos. The gap is the percent difference in each group's mean relative to the overall 2025 mean. The ratio divides each group's mean by the lowest mean within the same dimension. The sample excludes employed pension recipients. Sector rows omit extraterritorial organizations.}
\end{frame}

\begin{frame}{Paper Table 2 (2/3)}
\begin{minipage}{\textwidth}\small\textbf{Mean real hourly labor income gaps by worker group, 2025}\end{minipage}\par\vspace{0.15cm}
\centering\begingroup\fontsize{8.5}{10}\selectfont
\setlength{\tabcolsep}{4pt}\renewcommand{\arraystretch}{1.06}
\begin{adjustbox}{max width=\textwidth,max totalheight=0.64\textheight}
% Generated from Paper/sections/descriptive_wage_gaps_2025_table.tex
\begin{tabular}{lp{0.34\textwidth}rrrr}
\toprule
    Dimension & Group & Employment share & Mean income & Gap & Ratio \\
    \midrule

    Activity & A - Agriculture & 13.8\% & 6,036 & -41.4\% & 1.00 \\
& T - Household employers & 3.1\% & 6,454 & -37.3\% & 1.07 \\
& I - Accommodation and food & 7.6\% & 6,930 & -32.7\% & 1.15 \\
& H - Transport & 7.8\% & 7,775 & -24.5\% & 1.29 \\
& F - Construction & 6.9\% & 8,406 & -18.3\% & 1.39 \\
& G - Commerce & 17.0\% & 8,535 & -17.1\% & 1.41 \\
& R/S - Arts and other services & 5.4\% & 8,765 & -14.9\% & 1.45 \\
& C - Manufacturing & 11.0\% & 9,591 & -6.8\% & 1.59\\
\bottomrule\end{tabular}

\end{adjustbox}\endgroup\par\vspace{0.10cm}
{\fontsize{6.5}{7.5}\selectfont\raggedright Employment shares are computed within each dimension. Mean income is real hourly labor income in constant 2025 pesos. The gap is the percent difference in each group's mean relative to the overall 2025 mean. The ratio divides each group's mean by the lowest mean within the same dimension. The sample excludes employed pension recipients. Sector rows omit extraterritorial organizations.\par}
\note{Employment shares are computed within each dimension. Mean income is real hourly labor income in constant 2025 pesos. The gap is the percent difference in each group's mean relative to the overall 2025 mean. The ratio divides each group's mean by the lowest mean within the same dimension. The sample excludes employed pension recipients. Sector rows omit extraterritorial organizations.}
\end{frame}

\begin{frame}{Paper Table 2 (3/3)}
\begin{minipage}{\textwidth}\small\textbf{Mean real hourly labor income gaps by worker group, 2025}\end{minipage}\par\vspace{0.15cm}
\centering\begingroup\fontsize{8.5}{10}\selectfont
\setlength{\tabcolsep}{4pt}\renewcommand{\arraystretch}{1.06}
\begin{adjustbox}{max width=\textwidth,max totalheight=0.64\textheight}
% Generated from Paper/sections/descriptive_wage_gaps_2025_table.tex
\begin{tabular}{lp{0.34\textwidth}rrrr}
\toprule
    Dimension & Group & Employment share & Mean income & Gap & Ratio \\
    \midrule
Activity & D/E - Utilities and water & 1.3\% & 11,102 & 9.9\% & 1.89 \\
& B - Mining & 1.3\% & 11,136 & 10.2\% & 1.90 \\
& L/M/N - Business services & 9.0\% & 13,176 & 30.4\% & 2.24 \\
& Q - Health & 4.4\% & 14,376 & 42.3\% & 2.45 \\
& P - Education & 4.2\% & 19,335 & 91.4\% & 3.29 \\
& K - Finance and insurance & 1.9\% & 21,221 & 110.0\% & 3.62 \\
& O - Public administration & 3.6\% & 20,955 & 107.4\% & 3.57 \\
& J - Information and communication & 1.7\% & 21,608 & 113.9\% & 3.68\\
\bottomrule\end{tabular}

\end{adjustbox}\endgroup\par\vspace{0.10cm}
{\fontsize{6.5}{7.5}\selectfont\raggedright Employment shares are computed within each dimension. Mean income is real hourly labor income in constant 2025 pesos. The gap is the percent difference in each group's mean relative to the overall 2025 mean. The ratio divides each group's mean by the lowest mean within the same dimension. The sample excludes employed pension recipients. Sector rows omit extraterritorial organizations.\par}
\note{Employment shares are computed within each dimension. Mean income is real hourly labor income in constant 2025 pesos. The gap is the percent difference in each group's mean relative to the overall 2025 mean. The ratio divides each group's mean by the lowest mean within the same dimension. The sample excludes employed pension recipients. Sector rows omit extraterritorial organizations.}
\end{frame}

\begin{frame}{Paper Table 3}
\begin{minipage}{\textwidth}\small\textbf{Worker characteristics by firm size, 2025}\end{minipage}\par\vspace{0.15cm}
\centering\begingroup\fontsize{8.5}{10}\selectfont
\setlength{\tabcolsep}{4pt}\renewcommand{\arraystretch}{1.06}
\begin{adjustbox}{max width=\textwidth,max totalheight=0.64\textheight}
% Generated from Paper/sections/descriptive_worker_characteristics_2025_table.tex
\begin{tabular}{lrrrr}
\toprule
    Firm size & Women & Formal & Higher education & Mean age \\
    \midrule
Solo & 41.3\% & 8.3\% & 16.7\% & 44.5 \\
2-3 & 34.6\% & 13.1\% & 17.0\% & 41.3 \\
4-5 & 36.5\% & 24.2\% & 20.2\% & 37.6 \\
6-10 & 40.5\% & 46.6\% & 29.9\% & 37.0 \\
11-19 & 40.4\% & 64.4\% & 38.8\% & 37.2 \\
20-30 & 42.7\% & 81.2\% & 45.4\% & 37.0 \\
31-50 & 41.2\% & 90.4\% & 51.0\% & 37.5 \\
51-100 & 43.7\% & 94.3\% & 53.4\% & 37.5 \\
101+ & 45.8\% & 98.0\% & 62.9\% & 38.0\\
\bottomrule\end{tabular}

\end{adjustbox}\endgroup\par\vspace{0.10cm}
{\fontsize{6.5}{7.5}\selectfont\raggedright Higher education denotes education beyond secondary school.\par}
\note{Higher education denotes education beyond secondary school.}
\end{frame}

\begin{frame}{Paper Table 4}
\begin{minipage}{\textwidth}\small\textbf{Profile of formal workers by firm-size group, 2025}\end{minipage}\par\vspace{0.15cm}
\centering\begingroup\fontsize{8.5}{10}\selectfont
\setlength{\tabcolsep}{4pt}\renewcommand{\arraystretch}{1.06}
\begin{adjustbox}{max width=\textwidth,max totalheight=0.64\textheight}
% Generated from Paper/sections/descriptive_formal_worker_profile_2025_table.tex
\begin{tabular}{lrrr}
\toprule
    Characteristic & Formal solo & Formal 2--100 & Formal 101+ \\
    \midrule
\multicolumn{4}{l}{\textit{Employment shares}} \\
Share of formal workers & 6.4\% & 38.7\% & 54.9\% \\
\midrule
    \multicolumn{4}{l}{\textit{Remuneration}} \\
Mean real hourly income (2025 pesos) & 15,678 & 12,443 & 16,633 \\
\midrule
    \multicolumn{4}{l}{\textit{Personal characteristics}} \\
Mean age & 45.0 & 38.3 & 37.9 \\
Women & 43.3\% & 41.9\% & 45.7\% \\
Higher education & 50.8\% & 48.2\% & 63.4\% \\
\midrule
    \multicolumn{4}{l}{\textit{Status in employment}} \\
Own-account workers & 82.1\% & 7.7\% & 9.6\% \\
Domestic workers & 16.2\% & 0.9\% & 0.0\%\\
\bottomrule\end{tabular}

\end{adjustbox}\endgroup\par\vspace{0.10cm}
{\fontsize{6.5}{7.5}\selectfont\raggedright Employment shares use all formal workers as the denominator. Other percentages are within each firm-size group. The status-in-employment categories shown are not an exhaustive list. Higher education denotes education beyond secondary school.\par}
\note{Employment shares use all formal workers as the denominator. Other percentages are within each firm-size group. The status-in-employment categories shown are not an exhaustive list. Higher education denotes education beyond secondary school.}
\end{frame}

\begin{frame}{Paper Table 5}
\begin{minipage}{\textwidth}\small\textbf{Firm-size wage premiums with sequential controls}\end{minipage}\par\vspace{0.15cm}
\centering\begingroup\fontsize{8.5}{10}\selectfont
\setlength{\tabcolsep}{4pt}\renewcommand{\arraystretch}{1.06}
\begin{adjustbox}{max width=\textwidth,max totalheight=0.64\textheight}
% Generated from Paper/sections/regression_local_geography_formality_robustness_table.tex
\begin{tabular}{lcccc}
\toprule
    & \multicolumn{4}{c}{Dependent variable: log real hourly labor income} \\
    \cmidrule(lr){2-5}
    & (1) & (2) & (3) & \textbf{(4)} \\
    \midrule
\multicolumn{5}{l}{\textit{Informal workers}} \\
Firm size: 2-3 & 0.197*** \sereg{0.010} & 0.176*** \sereg{0.010} & 0.176*** \sereg{0.010} & 0.113*** \sereg{0.013} \\
Firm size: 4-5 & 0.343*** \sereg{0.012} & 0.311*** \sereg{0.011} & 0.308*** \sereg{0.011} & 0.255*** \sereg{0.016} \\
Firm size: 6-10 & 0.392*** \sereg{0.013} & 0.352*** \sereg{0.012} & 0.349*** \sereg{0.012} & 0.324*** \sereg{0.017} \\
Firm size: 11-19 & 0.408*** \sereg{0.016} & 0.368*** \sereg{0.014} & 0.364*** \sereg{0.014} & 0.354*** \sereg{0.018} \\
Firm size: 20-30 & 0.401*** \sereg{0.016} & 0.359*** \sereg{0.013} & 0.354*** \sereg{0.014} & 0.362*** \sereg{0.017} \\
Firm size: 31-50 & 0.399*** \sereg{0.018} & 0.359*** \sereg{0.018} & 0.351*** \sereg{0.018} & 0.363*** \sereg{0.020} \\
Firm size: 51-100 & 0.409*** \sereg{0.020} & 0.369*** \sereg{0.019} & 0.359*** \sereg{0.019} & 0.372*** \sereg{0.020} \\
Firm size: 101+ & 0.336*** \sereg{0.020} & 0.304*** \sereg{0.020} & 0.288*** \sereg{0.019} & 0.316*** \sereg{0.019} \\
\midrule
    Worker controls & Yes & Yes & Yes & Yes \\
Occupational-position controls & No & No & No & Yes \\
Fixed effects & $s\times t$ & $s\times t$, $d\times t$ & $s\times d\times t$ & $s\times d\times t$ \\
Observations & 5,268,514 & 5,268,514 & 5,268,375 & 5,268,375 \\
$R^2$ & 0.388 & 0.406 & 0.416 & 0.431\\
\bottomrule\end{tabular}

\end{adjustbox}\endgroup\par\vspace{0.10cm}
{\fontsize{6.5}{7.5}\selectfont\raggedright The omitted category is solo workers. Column (4) corresponds to Equation~(2). Parentheses contain sector-clustered standard errors. * $p<0.10$, ** $p<0.05$, *** $p<0.01$.\par}
\note{The omitted category is solo workers. Column (4) corresponds to Equation~(2). Parentheses contain sector-clustered standard errors. * $p<0.10$, ** $p<0.05$, *** $p<0.01$.}
\end{frame}

\begin{frame}{Paper Table 6}
\begin{minipage}{\textwidth}\small\textbf{Common-premium estimates with geographic and status-in-employment controls}\end{minipage}\par\vspace{0.15cm}
\centering\begingroup\fontsize{8.5}{10}\selectfont
\setlength{\tabcolsep}{4pt}\renewcommand{\arraystretch}{1.06}
\begin{adjustbox}{max width=\textwidth,max totalheight=0.64\textheight}
% Generated from Paper/sections/regression_local_geography_formality_robustness_101_reference_table.tex
\begin{tabular}{lcccc}
\toprule
    & \multicolumn{4}{c}{Dependent variable: log real hourly labor income} \\
    \cmidrule(lr){2-5}
    & (1) & (2) & (3) & \textbf{(4)} \\
    \midrule
\multicolumn{5}{l}{\textit{Informal workers}} \\
Firm size: Solo & -0.336*** \sereg{0.020} & -0.304*** \sereg{0.020} & -0.288*** \sereg{0.019} & -0.316*** \sereg{0.019} \\
Firm size: 2-3 & -0.139*** \sereg{0.017} & -0.127*** \sereg{0.017} & -0.112*** \sereg{0.016} & -0.202*** \sereg{0.016} \\
Firm size: 4-5 & 0.007 \sereg{0.016} & 0.007 \sereg{0.016} & 0.020 \sereg{0.016} & -0.061*** \sereg{0.017} \\
Firm size: 6-10 & 0.056*** \sereg{0.016} & 0.048*** \sereg{0.016} & 0.060*** \sereg{0.016} & 0.008 \sereg{0.017} \\
Firm size: 11-19 & 0.072*** \sereg{0.016} & 0.065*** \sereg{0.016} & 0.075*** \sereg{0.016} & 0.038** \sereg{0.017} \\
Firm size: 20-30 & 0.066*** \sereg{0.016} & 0.055*** \sereg{0.016} & 0.066*** \sereg{0.016} & 0.046*** \sereg{0.017} \\
Firm size: 31-50 & 0.064*** \sereg{0.016} & 0.056*** \sereg{0.015} & 0.063*** \sereg{0.015} & 0.047*** \sereg{0.016} \\
Firm size: 51-100 & 0.073*** \sereg{0.016} & 0.065*** \sereg{0.015} & 0.071*** \sereg{0.015} & 0.057*** \sereg{0.015} \\
\midrule
    Worker controls & Yes & Yes & Yes & Yes \\
Occupational-position controls & No & No & No & Yes \\
Fixed effects & $s\times t$ & $s\times t$, $d\times t$ & $s\times d\times t$ & $s\times d\times t$ \\
Observations & 5,268,514 & 5,268,514 & 5,268,375 & 5,268,375 \\
$R^2$ & 0.388 & 0.406 & 0.416 & 0.431\\
\bottomrule\end{tabular}

\end{adjustbox}\endgroup\par\vspace{0.10cm}
{\fontsize{6.5}{7.5}\selectfont\raggedright Common-premium model with solo as the reference. All columns include worker characteristics and formality. FE: (1) $s\times t$, (2) $s\times t$ and $d\times t$, (3)--(4) $s\times d\times t$. Column (4) adds status in employment. SEs cluster by sector-department. SEs in parentheses. * $p<0.10$, ** $p<0.05$, *** $p<0.01$.\par}
\note{The omitted category is solo workers. Column (4) corresponds to Equation~(3). $s$, $d$, and $t$ denote sector, department, and year. Parentheses contain sector-department-clustered standard errors. * $p<0.10$, ** $p<0.05$, *** $p<0.01$.}
\end{frame}

\begin{frame}{Paper Table 7 (1/2)}
\begin{minipage}{\textwidth}\small\textbf{Firm-size wage premiums by formality}\end{minipage}\par\vspace{0.15cm}
\centering\begingroup\fontsize{8.5}{10}\selectfont
\setlength{\tabcolsep}{4pt}\renewcommand{\arraystretch}{1.06}
\begin{adjustbox}{max width=\textwidth,max totalheight=0.64\textheight}
% Generated from Paper/sections/regression_local_geography_formality_robustness_11_19_reference_table.tex
\begin{tabular}{lcccc}
\toprule
    & \multicolumn{4}{c}{Dependent variable: log real hourly labor income} \\
    \cmidrule(lr){2-5}
    & (1) & (2) & (3) & \textbf{(4)} \\
    \midrule
\multicolumn{5}{l}{\textit{Informal workers}} \\
Firm size: Solo & -0.408*** \sereg{0.016} & -0.368*** \sereg{0.014} & -0.364*** \sereg{0.014} & -0.354*** \sereg{0.018} \\
Firm size: 2-3 & -0.211*** \sereg{0.013} & -0.192*** \sereg{0.011} & -0.187*** \sereg{0.011} & -0.240*** \sereg{0.011} \\
Firm size: 4-5 & -0.065*** \sereg{0.009} & -0.058*** \sereg{0.009} & -0.056*** \sereg{0.009} & -0.099*** \sereg{0.007} \\
Firm size: 6-10 & -0.016*** \sereg{0.006} & -0.016*** \sereg{0.006} & -0.015** \sereg{0.006} & -0.030*** \sereg{0.006} \\
Firm size: 20-30 & -0.007 \sereg{0.008} & -0.010 \sereg{0.007} & -0.009 \sereg{0.007} & 0.008 \sereg{0.007} \\
Firm size: 31-50 & -0.009 \sereg{0.011} & -0.009 \sereg{0.011} & -0.012 \sereg{0.011} & 0.009 \sereg{0.010} \\
Firm size: 51-100 & 0.001 \sereg{0.014} & 0.000 \sereg{0.013} & -0.004 \sereg{0.014} & 0.019 \sereg{0.013} \\
Firm size: 101+ & -0.072*** \sereg{0.016} & -0.065*** \sereg{0.016} & -0.075*** \sereg{0.016} & -0.038** \sereg{0.017} \\
\midrule
    Worker controls & Yes & Yes & Yes & Yes \\
Occupational-position controls & No & No & No & Yes \\
Fixed effects & $s\times t$ & $s\times t$, $d\times t$ & $s\times d\times t$ & $s\times d\times t$ \\
Observations & 5,268,514 & 5,268,514 & 5,268,375 & 5,268,375 \\
$R^2$ & 0.388 & 0.406 & 0.416 & 0.431\\
\bottomrule\end{tabular}

\end{adjustbox}\endgroup\par\vspace{0.10cm}
{\fontsize{6.5}{7.5}\selectfont\raggedright Within-formality log contrasts relative to solo workers. All columns include worker controls and GEIH weights. FE: (1) $s\times t$, (2) $s\times t$ and $d\times t$, (3)--(4) $s\times d\times t$. Column (4) adds status in employment and is preferred. SEs cluster by sector-department. SEs in parentheses. * $p<0.10$, ** $p<0.05$, *** $p<0.01$.\par}
\note{Within-formality log-wage contrasts relative to solo workers. Each column comes from one interacted regression. Column (4) corresponds to Equation~(4). $s$, $d$, and $t$ denote sector, department, and year. Parentheses contain sector-department-clustered standard errors. * $p<0.10$, ** $p<0.05$, *** $p<0.01$.}
\end{frame}

\begin{frame}{Paper Table 7 (2/2)}
\begin{minipage}{\textwidth}\small\textbf{Firm-size wage premiums by formality}\end{minipage}\par\vspace{0.15cm}
\centering\begingroup\fontsize{8.5}{10}\selectfont
\setlength{\tabcolsep}{4pt}\renewcommand{\arraystretch}{1.06}
\begin{adjustbox}{max width=\textwidth,max totalheight=0.64\textheight}
% Generated from Paper/sections/regression_local_geography_formality_robustness_11_19_reference_table.tex
\begin{tabular}{lcccc}
\toprule
    & \multicolumn{4}{c}{Dependent variable: log real hourly labor income} \\
    \cmidrule(lr){2-5}
    & (1) & (2) & (3) & \textbf{(4)} \\
    \midrule
\multicolumn{5}{l}{\textit{Formal workers}} \\
Firm size: Solo & -0.035 \sereg{0.034} & -0.038 \sereg{0.035} & -0.033 \sereg{0.034} & -0.039 \sereg{0.032} \\
Firm size: 2-3 & -0.013 \sereg{0.014} & -0.014 \sereg{0.014} & -0.013 \sereg{0.013} & -0.112*** \sereg{0.014} \\
Firm size: 4-5 & -0.010 \sereg{0.009} & -0.011 \sereg{0.008} & -0.010 \sereg{0.008} & -0.062*** \sereg{0.008} \\
Firm size: 6-10 & -0.012* \sereg{0.007} & -0.015** \sereg{0.007} & -0.014** \sereg{0.007} & -0.027*** \sereg{0.006} \\
Firm size: 20-30 & -0.011** \sereg{0.005} & -0.013*** \sereg{0.004} & -0.012*** \sereg{0.004} & -0.000 \sereg{0.004} \\
Firm size: 31-50 & 0.007 \sereg{0.005} & 0.008 \sereg{0.005} & 0.007 \sereg{0.005} & 0.020*** \sereg{0.005} \\
Firm size: 51-100 & 0.034*** \sereg{0.007} & 0.033*** \sereg{0.007} & 0.033*** \sereg{0.007} & 0.046*** \sereg{0.007} \\
Firm size: 101+ & 0.119*** \sereg{0.015} & 0.123*** \sereg{0.016} & 0.114*** \sereg{0.016} & 0.113*** \sereg{0.014} \\
\midrule
    Worker controls & Yes & Yes & Yes & Yes \\
Occupational-position controls & No & No & No & Yes \\
Fixed effects & $s\times t$ & $s\times t$, $d\times t$ & $s\times d\times t$ & $s\times d\times t$ \\
Observations & 5,268,514 & 5,268,514 & 5,268,375 & 5,268,375 \\
$R^2$ & 0.388 & 0.406 & 0.416 & 0.431\\
\bottomrule\end{tabular}

\end{adjustbox}\endgroup\par\vspace{0.10cm}
{\fontsize{6.5}{7.5}\selectfont\raggedright Within-formality log contrasts relative to solo workers. All columns include worker controls and GEIH weights. FE: (1) $s\times t$, (2) $s\times t$ and $d\times t$, (3)--(4) $s\times d\times t$. Column (4) adds status in employment and is preferred. SEs cluster by sector-department. SEs in parentheses. * $p<0.10$, ** $p<0.05$, *** $p<0.01$.\par}
\note{Within-formality log-wage contrasts relative to solo workers. Each column comes from one interacted regression. Column (4) corresponds to Equation~(4). $s$, $d$, and $t$ denote sector, department, and year. Parentheses contain sector-department-clustered standard errors. * $p<0.10$, ** $p<0.05$, *** $p<0.01$.}
\end{frame}

\begin{frame}{Paper Table 8}
\begin{minipage}{\textwidth}\small\textbf{Formal-informal wage gaps in the preferred specification}\end{minipage}\par\vspace{0.15cm}
\centering\begingroup\fontsize{8.5}{10}\selectfont
\setlength{\tabcolsep}{4pt}\renewcommand{\arraystretch}{1.06}
\begin{adjustbox}{max width=\textwidth,max totalheight=0.64\textheight}
% Generated from Paper/sections/regression_preferred_formality_gap_table.tex
\begin{tabular}{lccc}
\toprule
    Firm size & Log-wage gap & Premium & $p$-value \\
    \midrule
Solo & 0.468*** & 59.6\% & $<0.001$ \\
2-3 & 0.281*** & 32.4\% & $<0.001$ \\
4-5 & 0.188*** & 20.7\% & $<0.001$ \\
6-10 & 0.157*** & 17.0\% & $<0.001$ \\
11-19 & 0.155*** & 16.7\% & $<0.001$ \\
20-30 & 0.146*** & 15.8\% & $<0.001$ \\
31-50 & 0.169*** & 18.4\% & $<0.001$ \\
51-100 & 0.182*** & 19.9\% & $<0.001$ \\
101+ & 0.312*** & 36.6\% & $<0.001$\\
\bottomrule\end{tabular}

\end{adjustbox}\endgroup\par\vspace{0.10cm}
{\fontsize{6.5}{7.5}\selectfont\raggedright Formal-minus-informal contrasts from Equation~(4) at each firm size. Percentage premiums use informal workers as the denominator. Two-sided tests use sector-department-clustered standard errors. * $p<0.10$, ** $p<0.05$, *** $p<0.01$.\par}
\note{Formal-minus-informal contrasts from Equation~(4) at each firm size. Percentage premiums use informal workers as the denominator. Two-sided tests use sector-department-clustered standard errors. * $p<0.10$, ** $p<0.05$, *** $p<0.01$.}
\end{frame}

\begin{frame}{Paper Table 9 (1/2)}
\begin{minipage}{\textwidth}\small\textbf{Firm-size premiums in the full and salaried-only samples}\end{minipage}\par\vspace{0.15cm}
\centering\begingroup\fontsize{8.5}{10}\selectfont
\setlength{\tabcolsep}{4pt}\renewcommand{\arraystretch}{1.06}
\begin{adjustbox}{max width=\textwidth,max totalheight=0.64\textheight}
% Generated from Paper/sections/regression_salaried_robustness_table.tex
\begin{tabular}{lrrrr}
\toprule
    & \multicolumn{2}{c}{All workers} & \multicolumn{2}{c}{Salaried workers} \\
    \cmidrule(lr){2-3}\cmidrule(lr){4-5}
    Firm size & Informal & Formal & Informal & Formal \\
    \midrule
Solo & -0.354*** \sereg{0.018} & -0.039 \sereg{0.032} & -0.221*** \sereg{0.016} & -0.063*** \sereg{0.022} \\
2--3 & -0.240*** \sereg{0.011} & -0.112*** \sereg{0.014} & -0.178*** \sereg{0.010} & -0.080*** \sereg{0.013} \\
4--5 & -0.099*** \sereg{0.007} & -0.062*** \sereg{0.008} & -0.081*** \sereg{0.007} & -0.056*** \sereg{0.007} \\
6--10 & -0.030*** \sereg{0.006} & -0.027*** \sereg{0.006} & -0.029*** \sereg{0.005} & -0.021*** \sereg{0.005} \\
11--19 & Reference & Reference & Reference & Reference \\
20--30 & 0.008 \sereg{0.007} & -0.000 \sereg{0.004} & 0.007 \sereg{0.007} & 0.009*** \sereg{0.003} \\
31--50 & 0.009 \sereg{0.010} & 0.020*** \sereg{0.005} & 0.009 \sereg{0.010} & 0.032*** \sereg{0.004} \\
51--100 & 0.019 \sereg{0.013} & 0.046*** \sereg{0.007} & 0.023* \sereg{0.013} & 0.060*** \sereg{0.006} \\
101+ & -0.038** \sereg{0.017} & 0.113*** \sereg{0.014} & -0.033** \sereg{0.016} & 0.121*** \sereg{0.012}\\
\bottomrule\end{tabular}

\end{adjustbox}\endgroup\par\vspace{0.10cm}
{\fontsize{6.5}{7.5}\selectfont\raggedright Main rows are log contrasts relative to 11--19 workers within formality. Summary premiums instead use solo work. Same preferred controls and weights. Salaried workers include private, government, and domestic employees and day laborers. SEs cluster by sector-department (559 full-sample, 555 salaried clusters). Equality tests use the full covariance matrix. SEs in parentheses. * $p<0.10$, ** $p<0.05$, *** $p<0.01$.\par}
\note{Main rows report within-formality log-wage contrasts relative to 11--19 workers. Summary premiums use solo work as the reference and are in percent. Each column pair estimates Equation~(4) on the indicated sample. The restricted sample excludes 66 additional singleton cells. Standard errors in parentheses use sector-department clustering, with $G=559$ clusters in the full sample and $G=555$ among employees. Equality tests compare formal and informal log contrasts relative to solo work, individually for 101+ (two-sided $t$ test) and jointly across all sizes (Wald $F$ test with eight numerator degrees of freedom). Both use $G-1$ denominator degrees of freedom. * $p<0.10$, ** $p<0.05$, *** $p<0.01$.}
\end{frame}

\begin{frame}{Paper Table 9 (2/2)}
\begin{minipage}{\textwidth}\small\textbf{Firm-size premiums in the full and salaried-only samples}\end{minipage}\par\vspace{0.15cm}
\centering\begingroup\fontsize{8.5}{10}\selectfont
\setlength{\tabcolsep}{4pt}\renewcommand{\arraystretch}{1.06}
\begin{adjustbox}{max width=\textwidth,max totalheight=0.64\textheight}
% Generated from Paper/sections/regression_salaried_robustness_table.tex
\begin{tabular}{lrrrr}
\toprule
    & \multicolumn{2}{c}{All workers} & \multicolumn{2}{c}{Salaried workers} \\
    \cmidrule(lr){2-3}\cmidrule(lr){4-5}
    Firm size & Informal & Formal & Informal & Formal \\
    \midrule

    11--19/solo premium (\%) & 42.4 & 4.0 & 24.7 & 6.5 \\
101+/solo premium (\%) & 37.1 & 16.5 & 20.6 & 20.2 \\
Equal 101+/solo premiums: $p$ & \multicolumn{2}{c}{$<0.001$} & \multicolumn{2}{c}{0.897} \\
Equal size profiles: $p$ & \multicolumn{2}{c}{$<0.001$} & \multicolumn{2}{c}{$<0.001$} \\
Observations & \multicolumn{2}{c}{5,268,375} & \multicolumn{2}{c}{2,682,009}\\
\bottomrule\end{tabular}

\end{adjustbox}\endgroup\par\vspace{0.10cm}
{\fontsize{6.5}{7.5}\selectfont\raggedright Main rows are log contrasts relative to 11--19 workers within formality. Summary premiums instead use solo work. Same preferred controls and weights. Salaried workers include private, government, and domestic employees and day laborers. SEs cluster by sector-department (559 full-sample, 555 salaried clusters). Equality tests use the full covariance matrix. SEs in parentheses. * $p<0.10$, ** $p<0.05$, *** $p<0.01$.\par}
\note{Main rows report within-formality log-wage contrasts relative to 11--19 workers. Summary premiums use solo work as the reference and are in percent. Each column pair estimates Equation~(4) on the indicated sample. The restricted sample excludes 66 additional singleton cells. Standard errors in parentheses use sector-department clustering, with $G=559$ clusters in the full sample and $G=555$ among employees. Equality tests compare formal and informal log contrasts relative to solo work, individually for 101+ (two-sided $t$ test) and jointly across all sizes (Wald $F$ test with eight numerator degrees of freedom). Both use $G-1$ denominator degrees of freedom. * $p<0.10$, ** $p<0.05$, *** $p<0.01$.}
\end{frame}

\begin{frame}{Paper Table 10}
\begin{minipage}{\textwidth}\small\textbf{Changes in the 101+/solo premium in the full and salaried-only samples}\end{minipage}\par\vspace{0.15cm}
\centering\begingroup\fontsize{8.5}{10}\selectfont
\setlength{\tabcolsep}{4pt}\renewcommand{\arraystretch}{1.06}
\begin{adjustbox}{max width=\textwidth,max totalheight=0.64\textheight}
% Generated from Paper/sections/regression_preferred_endpoint_changes_table.tex
\begin{tabular}{lrrrr}
\toprule
    Firm size & 2008 (\%) & 2025 (\%) & Change (pp) & $p$-value \\
    \midrule
\multicolumn{5}{l}{\textit{Formal workers}} \\
2-3 & -11.9 & -10.9 & 1.0 & 0.784 \\
4-5 & -5.4 & -2.7 & 2.7 & 0.530 \\
6-10 & 0.1 & -0.7 & -0.8 & 0.865 \\
11-19 & 3.4 & 2.8 & -0.6 & 0.907 \\
20-30 & 2.1 & 1.0 & -1.1 & 0.820 \\
31-50 & 2.3 & 4.1 & 1.7 & 0.705 \\
51-100 & 3.7 & 6.2 & 2.5 & 0.613 \\
101+ & 12.6 & 13.5 & 0.9 & 0.861\\
\bottomrule\end{tabular}

\end{adjustbox}\endgroup\par\vspace{0.10cm}
{\fontsize{6.5}{7.5}\selectfont\raggedright Equation~(5) is estimated jointly for 2008 and 2025 on each sample. Premiums compare 101+ with solo workers of the same formality status. Changes are in percentage points. Two-sided tests concern the change in the underlying log-wage contrast, using the indicated clustering.\par}
\note{Equation~(5) is estimated jointly for 2008 and 2025 on each sample. Premiums compare 101+ with solo workers of the same formality status. Changes are in percentage points. Two-sided tests concern the change in the underlying log-wage contrast, using the indicated clustering.}
\end{frame}

\begin{frame}{Paper Table 11}
\begin{minipage}{\textwidth}\small\textbf{Income scenarios in the full and salaried-only samples, 2025}\end{minipage}\par\vspace{0.15cm}
\centering\begingroup\fontsize{8.5}{10}\selectfont
\setlength{\tabcolsep}{4pt}\renewcommand{\arraystretch}{1.06}
\begin{adjustbox}{max width=\textwidth,max totalheight=0.64\textheight}
% Generated from Paper/sections/income_salaried_comparison_table.tex
\begin{tabular}{@{}>{\raggedright\arraybackslash}p{0.36\textwidth}rrr@{}}
\toprule
    Scenario & \shortstack{Workers\\reassigned\\(millions)} & \shortstack{Reassigned\\mean income\\change (\%)} & \shortstack{Sample total\\income\\change (\%)} \\
    \midrule
\multicolumn{4}{@{}l}{\textit{A. All workers (22.56 million)}} \\
Informal to formal, same size & 12.43 & 42.1 & 13.5 \\
Formal to formal 101+ & 4.57 & 14.2 & 3.8 \\
Informal to informal 101+ & 12.31 & 24.0 & 7.6 \\
All workers to formal 101+ & 17.00 & 44.2 & 25.9 \\
\midrule
    \multicolumn{4}{@{}l}{\textit{B. Salaried workers (12.71 million)}} \\
Informal to formal, same size & 4.11 & 23.9 & 4.3 \\
Formal to formal 101+ & 3.57 & 12.8 & 3.6 \\
Informal to informal 101+ & 4.04 & 8.7 & 1.5 \\
All workers to formal 101+ & 7.68 & 26.5 & 12.3\\
\bottomrule\end{tabular}

\end{adjustbox}\endgroup\par\vspace{0.10cm}
{\fontsize{6.5}{7.5}\selectfont\raggedright Equation~(8) uses the pooled coefficients from Equation~(4) re-estimated on each sample. Income is predicted weekly labor income. Reassigned means exclude workers already at the destination. Sample totals include all workers represented in the corresponding panel, including those not reassigned.\par}
\note{Equation~(8) uses the pooled coefficients from Equation~(4) re-estimated on each sample. Income is predicted weekly labor income. Reassigned means exclude workers already at the destination. Sample totals include all workers represented in the corresponding panel, including those not reassigned.}
\end{frame}

\begin{frame}{Paper Table 12 (1/2)}
\begin{minipage}{\textwidth}\small\textbf{Changes in within-formality firm-size premiums, 2008--2025}\end{minipage}\par\vspace{0.15cm}
\centering\begingroup\fontsize{8.5}{10}\selectfont
\setlength{\tabcolsep}{4pt}\renewcommand{\arraystretch}{1.06}
\begin{adjustbox}{max width=\textwidth,max totalheight=0.64\textheight}
% Generated from Paper/sections/regression_preferred_endpoint_changes_table.tex
\begin{tabular}{lrrrr}
\toprule
    Firm size & 2008 (\%) & 2025 (\%) & Change (pp) & $p$-value \\
    \midrule
\multicolumn{5}{l}{\textit{Formal workers}} \\
2-3 & -7.7 & -11.2 & -3.5 & 0.265 \\
4-5 & -2.6 & -3.1 & -0.5 & 0.883 \\
6-10 & 4.1 & -1.8 & -5.8 & 0.117 \\
11-19 & 5.5 & 2.4 & -3.1 & 0.475 \\
20-30 & 5.5 & 0.9 & -4.7 & 0.266 \\
31-50 & 6.2 & 4.0 & -2.2 & 0.564 \\
51-100 & 5.7 & 6.3 & 0.5 & 0.899 \\
101+ & 16.6 & 13.5 & -3.2 & 0.467\\
\bottomrule\end{tabular}

\end{adjustbox}\endgroup\par\vspace{0.10cm}
{\fontsize{6.5}{7.5}\selectfont\raggedright Estimates from Equation~(5), with solo work as the reference within each formality-year group. Changes are in percentage points and calculated before rounding. Two-sided $p$-values test equality of the endpoint log-wage contrasts using joint estimation and sector-department clustering. Tests are pointwise, not adjusted for multiple comparisons.\par}
\note{Estimates from Equation~(5), with solo work as the reference within each formality-year group. Changes are in percentage points and calculated before rounding. Two-sided $p$-values test equality of the endpoint log-wage contrasts using joint estimation and sector-department clustering. Tests are pointwise, not adjusted for multiple comparisons.}
\end{frame}

\begin{frame}{Paper Table 12 (2/2)}
\begin{minipage}{\textwidth}\small\textbf{Changes in within-formality firm-size premiums, 2008--2025}\end{minipage}\par\vspace{0.15cm}
\centering\begingroup\fontsize{8.5}{10}\selectfont
\setlength{\tabcolsep}{4pt}\renewcommand{\arraystretch}{1.06}
\begin{adjustbox}{max width=\textwidth,max totalheight=0.64\textheight}
% Generated from Paper/sections/regression_preferred_endpoint_changes_table.tex
\begin{tabular}{lrrrr}
\toprule
    Firm size & 2008 (\%) & 2025 (\%) & Change (pp) & $p$-value \\
    \midrule
\multicolumn{5}{l}{\textit{Informal workers}} \\
2-3 & 7.5 & 18.2 & 10.7 & $<0.001$ \\
4-5 & 23.3 & 36.8 & 13.5 & 0.002 \\
6-10 & 29.4 & 47.2 & 17.8 & $<0.001$ \\
11-19 & 37.2 & 49.6 & 12.4 & 0.034 \\
20-30 & 44.5 & 44.5 & 0.0 & 0.997 \\
31-50 & 39.8 & 42.5 & 2.7 & 0.679 \\
51-100 & 43.6 & 47.3 & 3.7 & 0.639 \\
101+ & 33.2 & 57.1 & 23.8 & 0.002\\
\bottomrule\end{tabular}

\end{adjustbox}\endgroup\par\vspace{0.10cm}
{\fontsize{6.5}{7.5}\selectfont\raggedright Estimates from Equation~(5), with solo work as the reference within each formality-year group. Changes are in percentage points and calculated before rounding. Two-sided $p$-values test equality of the endpoint log-wage contrasts using joint estimation and sector-department clustering. Tests are pointwise, not adjusted for multiple comparisons.\par}
\note{Estimates from Equation~(5), with solo work as the reference within each formality-year group. Changes are in percentage points and calculated before rounding. Two-sided $p$-values test equality of the endpoint log-wage contrasts using joint estimation and sector-department clustering. Tests are pointwise, not adjusted for multiple comparisons.}
\end{frame}

\begin{frame}{Paper Table 13}
\begin{minipage}{\textwidth}\small\textbf{Approximate firm-size elasticities with preferred controls}\end{minipage}\par\vspace{0.15cm}
\centering\begingroup\fontsize{8.5}{10}\selectfont
\setlength{\tabcolsep}{4pt}\renewcommand{\arraystretch}{1.06}
\begin{adjustbox}{max width=\textwidth,max totalheight=0.64\textheight}
% Generated from Paper/sections/minimum_wage_bunching_2025_table.tex
\begin{tabular}{lrrrrrrr}
\toprule
    Firm size & Employment share & $<0.9$ & $0.9$--$1.1$ & $>1.1$ & P25 & Median & P75 \\
    \midrule
\multicolumn{8}{l}{\textit{Formal workers}} \\
Solo & 6.4\% & 21.6\% & 13.9\% & 64.5\% & 0.95 & 1.32 & 2.37 \\
2--3 & 4.0\% & 20.6\% & 28.4\% & 51.0\% & 0.94 & 1.12 & 1.58 \\
4--5 & 3.8\% & 13.9\% & 40.1\% & 46.0\% & 0.94 & 1.05 & 1.44 \\
6--10 & 6.0\% & 11.4\% & 42.7\% & 45.9\% & 0.94 & 1.07 & 1.40 \\
11--19 & 5.6\% & 9.0\% & 40.1\% & 50.9\% & 0.98 & 1.12 & 1.58 \\
20--30 & 5.9\% & 7.3\% & 40.9\% & 51.8\% & 0.98 & 1.12 & 1.58 \\
31--50 & 6.1\% & 7.4\% & 38.2\% & 54.4\% & 0.98 & 1.14 & 1.79 \\
51--100 & 7.3\% & 6.2\% & 36.9\% & 57.0\% & 0.98 & 1.17 & 1.90 \\
101+ & 54.9\% & 4.7\% & 28.0\% & 67.3\% & 1.02 & 1.38 & 2.59\\
\bottomrule\end{tabular}

\end{adjustbox}\endgroup\par\vspace{0.10cm}
{\fontsize{6.5}{7.5}\selectfont\raggedright Literature means are from Table 1 of Diegmann et al. (2026). All Colombian estimates use the preferred sample. Formal and informal slopes come from one regression, Equation~(6). The common-slope row imposes $\kappa=0$. Standard errors are clustered by sector-department. * $p<0.10$, ** $p<0.05$, *** $p<0.01$.\par}
\note{Literature means are from Table 1 of Diegmann et al. (2026). All Colombian estimates use the preferred sample. Formal and informal slopes come from one regression, Equation~(6). The common-slope row imposes $\kappa=0$. Standard errors are clustered by sector-department. * $p<0.10$, ** $p<0.05$, *** $p<0.01$.}
\end{frame}

\begin{frame}{Paper Table 14}
\begin{minipage}{\textwidth}\small\textbf{Illustrative changes in predicted labor income, 2025}\end{minipage}\par\vspace{0.15cm}
\centering\begingroup\fontsize{8.5}{10}\selectfont
\setlength{\tabcolsep}{4pt}\renewcommand{\arraystretch}{1.06}
\begin{adjustbox}{max width=\textwidth,max totalheight=0.64\textheight}
% Generated from Paper/sections/income_counterfactuals_2025_table.tex
\begin{tabular}{@{}>{\raggedright\arraybackslash}p{0.27\textwidth}lrrrr@{}}
\toprule
    Scenario & \shortstack[l]{Initial\\group} & \shortstack[r]{Workers\\reassigned\\(millions)} & \shortstack[r]{Group total\\(\%)} & \shortstack[r]{Reassigned\\mean (\%)} & \shortstack[r]{All-worker\\total (\%)} \\
    \midrule
Informal to formal, same size & Informal & 12.43 & 42.1 & 42.1 & 13.5 \\
Formal workers to formal 101+ & Formal & 4.57 & 5.5 & 14.2 & 3.8 \\
Informal workers to informal 101+ & Informal & 12.31 & 23.7 & 24.0 & 7.6 \\
\midrule
    All workers to formal 101+ & All & 17.00 & 25.9 & 44.2 & 25.9\\
\bottomrule\end{tabular}

\end{adjustbox}\endgroup\par\vspace{0.10cm}
{\fontsize{6.5}{7.5}\selectfont\raggedright Income columns report percentage changes in predicted weekly labor income from Equation~(8), using the pooled coefficients from Equation~(4). The 2025 sample represents 22.6 million workers. Group totals include everyone initially in the indicated group. Reassigned means use only workers whose size or formality changes. The all-worker total includes the entire 2025 sample. In the last row, both formal and informal workers are assigned to formal 101+ employment. Workers already there retain their baseline income.\par}
\note{Income columns report percentage changes in predicted weekly labor income from Equation~(8), using the pooled coefficients from Equation~(4). The 2025 sample represents 22.6 million workers. Group totals include everyone initially in the indicated group. Reassigned means use only workers whose size or formality changes. The all-worker total includes the entire 2025 sample. In the last row, both formal and informal workers are assigned to formal 101+ employment. Workers already there retain their baseline income.}
\end{frame}

\begin{frame}{Paper Table 15 (1/2)}
\begin{minipage}{\textwidth}\small\textbf{Distribution of hourly wages relative to the hourly minimum-wage benchmark, 2025}\end{minipage}\par\vspace{0.15cm}
\centering\begingroup\fontsize{8.5}{10}\selectfont
\setlength{\tabcolsep}{4pt}\renewcommand{\arraystretch}{1.06}
\begin{adjustbox}{max width=\textwidth,max totalheight=0.64\textheight}
% Generated from Paper/sections/descriptive_core_size_2008_2025_101_reference_table.tex
\begin{tabular}{lrrrrrrr}
\toprule
    & \multicolumn{2}{c}{Employment share} & \multicolumn{2}{c}{Mean hourly income} & \multicolumn{2}{c}{Wage ratio to 101+} & \multicolumn{1}{c}{Wage growth} \\
    \cmidrule(lr){2-3} \cmidrule(lr){4-5} \cmidrule(lr){6-7} \cmidrule(lr){8-8}
    Firm size & 2008 & 2025 & 2008 & 2025 & 2008 & 2025 & Annualized \\
    \midrule
Solo & 36.1\% & 34.7\% &  5,388 &  6,534 & 0.42 & 0.39 & 1.1\% \\
2--3 & 19.6\% & 13.8\% &  6,242 &  7,242 & 0.48 & 0.43 & 0.9\% \\
4--5 & 7.3\% & 7.0\% &  7,196 &  8,242 & 0.56 & 0.49 & 0.8\% \\
6--10 & 5.9\% & 5.8\% &  7,543 &  9,663 & 0.58 & 0.57 & 1.5\% \\
11--19 & 3.4\% & 3.9\% &  8,679 & 11,634 & 0.67 & 0.69 & 1.7\% \\
20--30 & 3.3\% & 3.3\% &  8,905 & 11,201 & 0.69 & 0.66 & 1.4\% \\
31--50 & 2.6\% & 3.0\% & 10,424 & 12,574 & 0.81 & 0.75 & 1.1\% \\
51--100 & 3.0\% & 3.5\% & 10,682 & 13,217 & 0.83 & 0.78 & 1.3\% \\
101+ & 18.8\% & 25.2\% & 12,943 & 16,859 & 1.00 & 1.00 & 1.6\% \\
\addlinespace
    All workers & 100.0\% & 100.0\% &  7,752 & 10,295 & 0.60 & 0.61 & 1.7\%\\
\bottomrule\end{tabular}

\end{adjustbox}\endgroup\par\vspace{0.10cm}
{\fontsize{6.5}{7.5}\selectfont\raggedright Wages are expressed relative to the hourly minimum-wage benchmark, computed as 1,423,500 pesos divided by $45\times52/12$ hours. This is an annual conversion benchmark, not a worker-specific legal compliance threshold. The columns $<0.9$, $0.9$--$1.1$, and $>1.1$ report the share within each size-formality group in each range. Employment shares are computed within each formality group. Quantiles refer to the wage-to-benchmark ratio.\par}
\note{Wages are expressed relative to the hourly minimum-wage benchmark, computed as 1,423,500 pesos divided by $45\times52/12$ hours. This is an annual conversion benchmark, not a worker-specific legal compliance threshold. The columns $<0.9$, $0.9$--$1.1$, and $>1.1$ report the share within each size-formality group in each range. Employment shares are computed within each formality group. Quantiles refer to the wage-to-benchmark ratio.}
\end{frame}

\begin{frame}{Paper Table 15 (2/2)}
\begin{minipage}{\textwidth}\small\textbf{Distribution of hourly wages relative to the hourly minimum-wage benchmark, 2025}\end{minipage}\par\vspace{0.15cm}
\centering\begingroup\fontsize{8.5}{10}\selectfont
\setlength{\tabcolsep}{4pt}\renewcommand{\arraystretch}{1.06}
\begin{adjustbox}{max width=\textwidth,max totalheight=0.64\textheight}
% Generated from Paper/sections/minimum_wage_bunching_2025_table.tex
\begin{tabular}{lrrrrrrr}
\toprule
    Firm size & Employment share & $<0.9$ & $0.9$--$1.1$ & $>1.1$ & P25 & Median & P75 \\
    \midrule
\multicolumn{8}{l}{\textit{Informal workers}} \\
Solo & 57.6\% & 74.3\% & 9.8\% & 15.9\% & 0.38 & 0.62 & 0.91 \\
2--3 & 21.7\% & 67.0\% & 14.9\% & 18.2\% & 0.47 & 0.71 & 0.99 \\
4--5 & 9.7\% & 61.1\% & 20.8\% & 18.1\% & 0.60 & 0.79 & 1.01 \\
6--10 & 5.6\% & 53.6\% & 21.7\% & 24.6\% & 0.63 & 0.86 & 1.09 \\
11--19 & 2.5\% & 49.3\% & 24.0\% & 26.6\% & 0.66 & 0.91 & 1.11 \\
20--30 & 1.1\% & 48.3\% & 24.6\% & 27.0\% & 0.66 & 0.91 & 1.12 \\
31--50 & 0.5\% & 47.6\% & 25.7\% & 26.7\% & 0.66 & 0.92 & 1.13 \\
51--100 & 0.4\% & 42.3\% & 21.9\% & 35.8\% & 0.68 & 0.98 & 1.28 \\
101+ & 0.9\% & 43.5\% & 19.5\% & 37.0\% & 0.66 & 0.95 & 1.37\\
\bottomrule\end{tabular}

\end{adjustbox}\endgroup\par\vspace{0.10cm}
{\fontsize{6.5}{7.5}\selectfont\raggedright Wages are expressed relative to the hourly minimum-wage benchmark, computed as 1,423,500 pesos divided by $45\times52/12$ hours. This is an annual conversion benchmark, not a worker-specific legal compliance threshold. The columns $<0.9$, $0.9$--$1.1$, and $>1.1$ report the share within each size-formality group in each range. Employment shares are computed within each formality group. Quantiles refer to the wage-to-benchmark ratio.\par}
\note{Wages are expressed relative to the hourly minimum-wage benchmark, computed as 1,423,500 pesos divided by $45\times52/12$ hours. This is an annual conversion benchmark, not a worker-specific legal compliance threshold. The columns $<0.9$, $0.9$--$1.1$, and $>1.1$ report the share within each size-formality group in each range. Employment shares are computed within each formality group. Quantiles refer to the wage-to-benchmark ratio.}
\end{frame}

\begin{frame}{Paper Table 16 (1/2)}
\begin{minipage}{\textwidth}\small\textbf{Distribution of monthly labor income relative to the monthly minimum wage, 2025}\end{minipage}\par\vspace{0.15cm}
\centering\begingroup\fontsize{8.5}{10}\selectfont
\setlength{\tabcolsep}{4pt}\renewcommand{\arraystretch}{1.06}
\begin{adjustbox}{max width=\textwidth,max totalheight=0.64\textheight}
% Generated from Paper/sections/descriptive_core_size_2008_2025_11_19_reference_table.tex
\begin{tabular}{lrrrrrrr}
\toprule
    & \multicolumn{2}{c}{Employment share} & \multicolumn{2}{c}{Mean hourly income} & \multicolumn{2}{c}{Wage ratio to 11--19} & \multicolumn{1}{c}{Wage growth} \\
    \cmidrule(lr){2-3} \cmidrule(lr){4-5} \cmidrule(lr){6-7} \cmidrule(lr){8-8}
    Firm size & 2008 & 2025 & 2008 & 2025 & 2008 & 2025 & Annualized \\
    \midrule
Solo & 36.1\% & 34.7\% &  5,388 &  6,534 & 0.62 & 0.56 & 1.1\% \\
2--3 & 19.6\% & 13.8\% &  6,242 &  7,242 & 0.72 & 0.62 & 0.9\% \\
4--5 & 7.3\% & 7.0\% &  7,196 &  8,242 & 0.83 & 0.71 & 0.8\% \\
6--10 & 5.9\% & 5.8\% &  7,543 &  9,663 & 0.87 & 0.83 & 1.5\% \\
11--19 & 3.4\% & 3.9\% &  8,679 & 11,634 & 1.00 & 1.00 & 1.7\% \\
20--30 & 3.3\% & 3.3\% &  8,905 & 11,201 & 1.03 & 0.96 & 1.4\% \\
31--50 & 2.6\% & 3.0\% & 10,424 & 12,574 & 1.20 & 1.08 & 1.1\% \\
51--100 & 3.0\% & 3.5\% & 10,682 & 13,217 & 1.23 & 1.14 & 1.3\% \\
101+ & 18.8\% & 25.2\% & 12,943 & 16,859 & 1.49 & 1.45 & 1.6\% \\
\addlinespace
    All workers & 100.0\% & 100.0\% &  7,752 & 10,295 & 0.89 & 0.88 & 1.7\%\\
\bottomrule\end{tabular}

\end{adjustbox}\endgroup\par\vspace{0.10cm}
{\fontsize{6.5}{7.5}\selectfont\raggedright Monthly labor income refers to the main job. The columns $<0.9$, $0.9$--$1.1$, and $>1.1$ report the share within each size-formality group in each range relative to the statutory monthly minimum wage. Employment shares are computed within each formality group. Quantiles refer to the income-to-minimum ratio.\par}
\note{Monthly labor income refers to the main job. The columns $<0.9$, $0.9$--$1.1$, and $>1.1$ report the share within each size-formality group in each range relative to the statutory monthly minimum wage. Employment shares are computed within each formality group. Quantiles refer to the income-to-minimum ratio.}
\end{frame}

\begin{frame}{Paper Table 16 (2/2)}
\begin{minipage}{\textwidth}\small\textbf{Distribution of monthly labor income relative to the monthly minimum wage, 2025}\end{minipage}\par\vspace{0.15cm}
\centering\begingroup\fontsize{8.5}{10}\selectfont
\setlength{\tabcolsep}{4pt}\renewcommand{\arraystretch}{1.06}
\begin{adjustbox}{max width=\textwidth,max totalheight=0.64\textheight}
% Generated from Paper/sections/minimum_wage_monthly_bunching_2025_table.tex
\begin{tabular}{lrrrrrrr}
\toprule
    Firm size & Employment share & $<0.9$ & $0.9$--$1.1$ & $>1.1$ & P25 & Median & P75 \\
    \midrule
\multicolumn{8}{l}{\textit{Informal workers}} \\
Solo & 57.6\% & 78.0\% & 10.1\% & 11.9\% & 0.28 & 0.53 & 0.84 \\
2--3 & 21.7\% & 65.5\% & 16.5\% & 18.0\% & 0.42 & 0.70 & 1.00 \\
4--5 & 9.7\% & 56.3\% & 22.5\% & 21.3\% & 0.56 & 0.84 & 1.05 \\
6--10 & 5.6\% & 50.9\% & 25.3\% & 23.7\% & 0.59 & 0.86 & 1.05 \\
11--19 & 2.5\% & 46.2\% & 26.0\% & 27.7\% & 0.63 & 0.91 & 1.12 \\
20--30 & 1.1\% & 46.3\% & 28.3\% & 25.4\% & 0.56 & 0.91 & 1.12 \\
31--50 & 0.5\% & 42.3\% & 28.4\% & 29.3\% & 0.70 & 0.98 & 1.14 \\
51--100 & 0.4\% & 41.9\% & 23.4\% & 34.7\% & 0.63 & 1.00 & 1.26 \\
101+ & 0.9\% & 45.9\% & 24.9\% & 29.2\% & 0.56 & 0.95 & 1.26\\
\bottomrule\end{tabular}

\end{adjustbox}\endgroup\par\vspace{0.10cm}
{\fontsize{6.5}{7.5}\selectfont\raggedright Monthly labor income refers to the main job. The columns $<0.9$, $0.9$--$1.1$, and $>1.1$ report the share within each size-formality group in each range relative to the statutory monthly minimum wage. Employment shares are computed within each formality group. Quantiles refer to the income-to-minimum ratio.\par}
\note{Monthly labor income refers to the main job. The columns $<0.9$, $0.9$--$1.1$, and $>1.1$ report the share within each size-formality group in each range relative to the statutory monthly minimum wage. Employment shares are computed within each formality group. Quantiles refer to the income-to-minimum ratio.}
\end{frame}

\begin{frame}{Paper Table 17 (1/2)}
\begin{minipage}{\textwidth}\small\textbf{Separating controls from clustering in the 101+/solo comparison}\end{minipage}\par\vspace{0.15cm}
\centering\begingroup\fontsize{8.5}{10}\selectfont
\setlength{\tabcolsep}{4pt}\renewcommand{\arraystretch}{1.06}
\begin{adjustbox}{max width=\textwidth,max totalheight=0.64\textheight}
% Generated from Paper/sections/regression_preferred_clustering_table.tex
\begin{tabular}{lrrrrr}
\toprule
    Controls & Premium (\%) & S.E. (SD) & $p$ (SD) & S.E. (S) & $p$ (S) \\
    \midrule
\multicolumn{6}{l}{\textit{Formal workers}} \\
Sector-year FE & 16.6 & 0.038 & $<0.001$ & 0.088 & 0.100 \\
Sector-department-year FE & 15.8 & 0.039 & $<0.001$ & 0.093 & 0.135 \\
+ Occupational position & 16.5 & 0.037 & $<0.001$ & 0.091 & 0.113\\
\bottomrule\end{tabular}

\end{adjustbox}\endgroup\par\vspace{0.10cm}
{\fontsize{6.5}{7.5}\selectfont\raggedright Specifications from columns (1), (3), and (4) of Table~7, estimated on the same 5,426,174 observations. The reference is solo work within each formality group. SD denotes sector-department clustering and S denotes sector clustering. Standard errors are in log points. Two-sided tests use 558 and 16 degrees of freedom, respectively.\par}
\note{Specifications from columns (1), (3), and (4) of Table~7, estimated on the same 5,426,174 observations. The reference is solo work within each formality group. SD denotes sector-department clustering and S denotes sector clustering. Standard errors are in log points. Two-sided tests use 558 and 16 degrees of freedom, respectively.}
\end{frame}

\begin{frame}{Paper Table 17 (2/2)}
\begin{minipage}{\textwidth}\small\textbf{Separating controls from clustering in the 101+/solo comparison}\end{minipage}\par\vspace{0.15cm}
\centering\begingroup\fontsize{8.5}{10}\selectfont
\setlength{\tabcolsep}{4pt}\renewcommand{\arraystretch}{1.06}
\begin{adjustbox}{max width=\textwidth,max totalheight=0.64\textheight}
% Generated from Paper/sections/regression_preferred_clustering_table.tex
\begin{tabular}{lrrrrr}
\toprule
    Controls & Premium (\%) & S.E. (SD) & $p$ (SD) & S.E. (S) & $p$ (S) \\
    \midrule
\multicolumn{6}{l}{\textit{Informal workers}} \\
Sector-year FE & 40.6 & 0.020 & $<0.001$ & 0.047 & $<0.001$ \\
Sector-department-year FE & 33.9 & 0.020 & $<0.001$ & 0.044 & $<0.001$ \\
+ Status in employment & 37.0 & 0.019 & $<0.001$ & 0.038 & $<0.001$\\
\bottomrule\end{tabular}

\end{adjustbox}\endgroup\par\vspace{0.10cm}
{\fontsize{6.5}{7.5}\selectfont\raggedright Specifications from columns (1), (3), and (4) of Table~7, estimated on the same 5,426,174 observations. The reference is solo work within each formality group. SD denotes sector-department clustering and S denotes sector clustering. Standard errors are in log points. Two-sided tests use 558 and 16 degrees of freedom, respectively.\par}
\note{Specifications from columns (1), (3), and (4) of Table~7, estimated on the same 5,426,174 observations. The reference is solo work within each formality group. SD denotes sector-department clustering and S denotes sector clustering. Standard errors are in log points. Two-sided tests use 558 and 16 degrees of freedom, respectively.}
\end{frame}

\begin{frame}{Paper Table 18 (1/2)}
\begin{minipage}{\textwidth}\small\textbf{Clustering sensitivity of preferred-model contrasts}\end{minipage}\par\vspace{0.15cm}
\centering\begingroup\fontsize{8.5}{10}\selectfont
\setlength{\tabcolsep}{4pt}\renewcommand{\arraystretch}{1.06}
\begin{adjustbox}{max width=\textwidth,max totalheight=0.64\textheight}
% Generated from Paper/sections/regression_preferred_extension_clustering_table.tex
\begin{tabular}{lrrrrr}
\toprule
    Workers & Estimate & S.E. (SD) & $p$ (SD) & S.E. (S) & $p$ (S) \\
    \midrule
\multicolumn{6}{l}{\textit{Change in the 101+/solo log-wage gap, 2008--2025}} \\
Formal & 0.008 & 0.045 & 0.861 & 0.047 & 0.869 \\
Informal & 0.152 & 0.060 & 0.011 & 0.084 & 0.090 \\
\midrule
    \multicolumn{6}{l}{\textit{Approximate firm-size elasticity}} \\
Formal & 0.040 & 0.005 & $<0.001$ & 0.012 & 0.003 \\
Informal & 0.097 & 0.004 & $<0.001$ & 0.008 & $<0.001$\\
\bottomrule\end{tabular}

\end{adjustbox}\endgroup\par\vspace{0.10cm}
{\fontsize{6.5}{7.5}\selectfont\raggedright Panels A--B use Equation~(4), C uses the joint endpoint model in Equation~(5), and D uses Equation~(6). Estimates in A--C are in log points, not percentages or percentage-point changes. SD denotes sector-department intersections, S sector, D department, and S+D two-way sector and department clustering. Two-sided $t$ tests use the number of clusters minus one degree of freedom, with the smaller cluster count for S+D.\par}
\note{Panels A--B use Equation~(4), C uses the joint endpoint model in Equation~(5), and D uses Equation~(6). Estimates in A--C are in log points, not percentages or percentage-point changes. SD denotes sector-department intersections, S sector, D department, and S+D two-way sector and department clustering. Two-sided $t$ tests use the number of clusters minus one degree of freedom, with the smaller cluster count for S+D.}
\end{frame}

\begin{frame}{Paper Table 18 (2/2)}
\begin{minipage}{\textwidth}\small\textbf{Clustering sensitivity of preferred-model contrasts}\end{minipage}\par\vspace{0.15cm}
\centering\begingroup\fontsize{8.5}{10}\selectfont
\setlength{\tabcolsep}{4pt}\renewcommand{\arraystretch}{1.06}
\begin{adjustbox}{max width=\textwidth,max totalheight=0.64\textheight}
% Generated from Paper/sections/regression_inference_clustering_table.tex
\begin{tabular}{@{}p{0.40\textwidth}rrrrr@{}}
\toprule
    Contrast & Estimate & $p$ (SD) & $p$ (S) & $p$ (D) & $p$ (S+D) \\
    \midrule

    \multicolumn{6}{l}{\textit{B. Formal-minus-informal log-wage gaps}} \\
Solo & 0.468 & $<0.001$ & $<0.001$ & $<0.001$ & $<0.001$ \\
11--19 & 0.155 & $<0.001$ & $<0.001$ & $<0.001$ & $<0.001$ \\
101+ & 0.312 & $<0.001$ & $<0.001$ & $<0.001$ & $<0.001$ \\
\midrule
    \multicolumn{6}{l}{\textit{C. Change in the 101+/solo log-wage gap, 2008--2025}} \\
Formal & -0.028 & 0.4672 & 0.5273 & 0.6621 & 0.6805 \\
Informal & 0.165 & 0.0018 & 0.0400 & 0.0608 & 0.1176 \\
\midrule
    \multicolumn{6}{l}{\textit{D. Approximate firm-size elasticities}} \\
Formal & 0.041 & $<0.001$ & 0.0021 & $<0.001$ & 0.0013 \\
Informal & 0.097 & $<0.001$ & $<0.001$ & $<0.001$ & $<0.001$ \\
Formal-minus-informal & -0.055 & $<0.001$ & $<0.001$ & $<0.001$ & $<0.001$\\
\bottomrule\end{tabular}

\end{adjustbox}\endgroup\par\vspace{0.10cm}
{\fontsize{6.5}{7.5}\selectfont\raggedright Panels A--B use Equation~(4), C uses the joint endpoint model in Equation~(5), and D uses Equation~(6). Estimates in A--C are in log points, not percentages or percentage-point changes. SD denotes sector-department intersections, S sector, D department, and S+D two-way sector and department clustering. Two-sided $t$ tests use the number of clusters minus one degree of freedom, with the smaller cluster count for S+D.\par}
\note{Panels A--B use Equation~(4), C uses the joint endpoint model in Equation~(5), and D uses Equation~(6). Estimates in A--C are in log points, not percentages or percentage-point changes. SD denotes sector-department intersections, S sector, D department, and S+D two-way sector and department clustering. Two-sided $t$ tests use the number of clusters minus one degree of freedom, with the smaller cluster count for S+D.}
\end{frame}

\begin{frame}{Paper Table 19}
\begin{minipage}{\textwidth}\small\textbf{Firm-size wage premium regressions: 101+ workers as the reference category}\end{minipage}\par\vspace{0.15cm}
\centering\begingroup\fontsize{8.5}{10}\selectfont
\setlength{\tabcolsep}{4pt}\renewcommand{\arraystretch}{1.06}
\begin{adjustbox}{max width=\textwidth,max totalheight=0.64\textheight}
% Generated from Paper/sections/regression_firm_size_table_101_reference.tex
\begin{tabular}{lcccc}
\toprule
    & \multicolumn{4}{c}{Dependent variable: log real hourly labor income} \\
    \cmidrule(lr){2-5}
    & (1) & (2) & (3) & (4) \\
    \midrule
Firm size: Solo & -1.017*** \sereg{0.102} & -0.817*** \sereg{0.067} & -0.598*** \sereg{0.061} & -0.359*** \sereg{0.059} \\
Firm size: 2-3 & -0.833*** \sereg{0.097} & -0.575*** \sereg{0.056} & -0.389*** \sereg{0.046} & -0.171*** \sereg{0.042} \\
Firm size: 4-5 & -0.654*** \sereg{0.099} & -0.409*** \sereg{0.052} & -0.237*** \sereg{0.040} & -0.052 \sereg{0.032} \\
Firm size: 6-10 & -0.518*** \sereg{0.093} & -0.317*** \sereg{0.051} & -0.176*** \sereg{0.037} & -0.043 \sereg{0.033} \\
Firm size: 11-19 & -0.391*** \sereg{0.082} & -0.242*** \sereg{0.055} & -0.141*** \sereg{0.041} & -0.057 \sereg{0.039} \\
Firm size: 20-30 & -0.330*** \sereg{0.070} & -0.214*** \sereg{0.051} & -0.135*** \sereg{0.039} & -0.088** \sereg{0.038} \\
Firm size: 31-50 & -0.258*** \sereg{0.060} & -0.167*** \sereg{0.045} & -0.110*** \sereg{0.034} & -0.088** \sereg{0.034} \\
Firm size: 51-100 & -0.198*** \sereg{0.052} & -0.119*** \sereg{0.036} & -0.079** \sereg{0.027} & -0.070** \sereg{0.028} \\
\midrule
    Gender, age, and education controls & No & No & Yes & Yes \\
Formality control & No & No & No & Yes \\
Sector-year fixed effects & No & Yes & Yes & Yes \\
Observations & 5,268,514 & 5,268,514 & 5,268,514 & 5,268,514 \\
$R^2$ & 0.213 & 0.271 & 0.376 & 0.386\\
\bottomrule\end{tabular}

\end{adjustbox}\endgroup\par\vspace{0.10cm}
{\fontsize{6.5}{7.5}\selectfont\raggedright Comparison specification. Specifications as in Table~5, with 101+ workers as the omitted category. Parentheses contain sector-clustered standard errors. * $p<0.10$, ** $p<0.05$, *** $p<0.01$.\par}
\note{Specifications as in Table~5, with 101+ workers as the omitted category. Parentheses contain sector-clustered standard errors. * $p<0.10$, ** $p<0.05$, *** $p<0.01$.}
\end{frame}

\begin{frame}{Paper Table 20}
\begin{minipage}{\textwidth}\small\textbf{Firm-size wage premium regressions: 11--19 workers as the reference category}\end{minipage}\par\vspace{0.15cm}
\centering\begingroup\fontsize{8.5}{10}\selectfont
\setlength{\tabcolsep}{4pt}\renewcommand{\arraystretch}{1.06}
\begin{adjustbox}{max width=\textwidth,max totalheight=0.64\textheight}
% Generated from Paper/sections/regression_firm_size_table_11_19_reference.tex
\begin{tabular}{lcccc}
\toprule
    & \multicolumn{4}{c}{Dependent variable: log real hourly labor income} \\
    \cmidrule(lr){2-5}
    & (1) & (2) & (3) & (4) \\
    \midrule
Firm size: Solo & -0.625*** \sereg{0.034} & -0.575*** \sereg{0.040} & -0.457*** \sereg{0.042} & -0.302*** \sereg{0.050} \\
Firm size: 2-3 & -0.442*** \sereg{0.029} & -0.333*** \sereg{0.033} & -0.248*** \sereg{0.028} & -0.114*** \sereg{0.030} \\
Firm size: 4-5 & -0.262*** \sereg{0.028} & -0.167*** \sereg{0.021} & -0.096*** \sereg{0.021} & 0.005 \sereg{0.020} \\
Firm size: 6-10 & -0.127*** \sereg{0.016} & -0.075*** \sereg{0.010} & -0.035*** \sereg{0.010} & 0.014 \sereg{0.010} \\
Firm size: 20-30 & 0.062*** \sereg{0.016} & 0.028*** \sereg{0.008} & 0.005 \sereg{0.006} & -0.031*** \sereg{0.005} \\
Firm size: 31-50 & 0.134*** \sereg{0.027} & 0.075*** \sereg{0.014} & 0.030*** \sereg{0.009} & -0.031*** \sereg{0.007} \\
Firm size: 51-100 & 0.193*** \sereg{0.035} & 0.123*** \sereg{0.024} & 0.061*** \sereg{0.017} & -0.014 \sereg{0.015} \\
Firm size: 101+ & 0.391*** \sereg{0.082} & 0.242*** \sereg{0.055} & 0.141*** \sereg{0.041} & 0.057 \sereg{0.039} \\
\midrule
    Gender, age, and education controls & No & No & Yes & Yes \\
Formality control & No & No & No & Yes \\
Sector-year fixed effects & No & Yes & Yes & Yes \\
Observations & 5,268,514 & 5,268,514 & 5,268,514 & 5,268,514 \\
$R^2$ & 0.213 & 0.271 & 0.376 & 0.386\\
\bottomrule\end{tabular}

\end{adjustbox}\endgroup\par\vspace{0.10cm}
{\fontsize{6.5}{7.5}\selectfont\raggedright Comparison specification. Specifications as in Table~5, with 11--19 workers as the omitted category. Parentheses contain sector-clustered standard errors. * $p<0.10$, ** $p<0.05$, *** $p<0.01$.\par}
\note{Specifications as in Table~5, with 11--19 workers as the omitted category. Parentheses contain sector-clustered standard errors. * $p<0.10$, ** $p<0.05$, *** $p<0.01$.}
\end{frame}

\begin{frame}{Paper Table 21}
\begin{minipage}{\textwidth}\small\textbf{Common-premium estimates with geographic and status-in-employment controls: 101+ workers as reference}\end{minipage}\par\vspace{0.15cm}
\centering\begingroup\fontsize{8.5}{10}\selectfont
\setlength{\tabcolsep}{4pt}\renewcommand{\arraystretch}{1.06}
\begin{adjustbox}{max width=\textwidth,max totalheight=0.64\textheight}
% Generated from Paper/sections/regression_local_geography_robustness_101_reference_table.tex
\begin{tabular}{lcccc}
\toprule
    & \multicolumn{4}{c}{Dependent variable: log real hourly labor income} \\
    \cmidrule(lr){2-5}
    & (1) & (2) & (3) & (4) \\
    \midrule
Firm size: Solo & -0.372*** \sereg{0.024} & -0.352*** \sereg{0.024} & -0.337*** \sereg{0.024} & -0.343*** \sereg{0.022} \\
Firm size: 2-3 & -0.179*** \sereg{0.018} & -0.178*** \sereg{0.018} & -0.165*** \sereg{0.017} & -0.242*** \sereg{0.016} \\
Firm size: 4-5 & -0.056*** \sereg{0.013} & -0.066*** \sereg{0.014} & -0.056*** \sereg{0.014} & -0.117*** \sereg{0.014} \\
Firm size: 6-10 & -0.049*** \sereg{0.012} & -0.063*** \sereg{0.013} & -0.053*** \sereg{0.013} & -0.078*** \sereg{0.012} \\
Firm size: 11-19 & -0.062*** \sereg{0.014} & -0.072*** \sereg{0.015} & -0.064*** \sereg{0.015} & -0.071*** \sereg{0.013} \\
Firm size: 20-30 & -0.090*** \sereg{0.014} & -0.100*** \sereg{0.015} & -0.091*** \sereg{0.014} & -0.082*** \sereg{0.012} \\
Firm size: 31-50 & -0.090*** \sereg{0.012} & -0.096*** \sereg{0.014} & -0.088*** \sereg{0.013} & -0.076*** \sereg{0.011} \\
Firm size: 51-100 & -0.072*** \sereg{0.012} & -0.079*** \sereg{0.013} & -0.070*** \sereg{0.012} & -0.058*** \sereg{0.011} \\
\midrule
    Worker and formality controls & Yes & Yes & Yes & Yes \\
Status-in-employment controls & No & No & No & Yes \\
Fixed effects & $s\times t$ & $s\times t$, $d\times t$ & $s\times d\times t$ & $s\times d\times t$ \\
Observations & 5,426,316 & 5,426,316 & 5,426,174 & 5,426,174 \\
$R^2$ & 0.391 & 0.410 & 0.420 & 0.435\\
\bottomrule\end{tabular}

\end{adjustbox}\endgroup\par\vspace{0.10cm}
{\fontsize{6.5}{7.5}\selectfont\raggedright Comparison specification. Common-premium model with 101+ as the reference. All columns include worker characteristics and formality. FE: (1) $s\times t$, (2) $s\times t$ and $d\times t$, (3)--(4) $s\times d\times t$. Column (4) adds status in employment. SEs cluster by sector-department. SEs in parentheses. * $p<0.10$, ** $p<0.05$, *** $p<0.01$.\par}
\note{Specifications as in Table~6, with 101+ workers as the omitted category. $s$, $d$, and $t$ denote sector, department, and year. Parentheses contain sector-department-clustered standard errors. * $p<0.10$, ** $p<0.05$, *** $p<0.01$.}
\end{frame}

\begin{frame}{Paper Table 22}
\begin{minipage}{\textwidth}\small\textbf{Common-premium estimates with geographic and status-in-employment controls: 11--19 workers as reference}\end{minipage}\par\vspace{0.15cm}
\centering\begingroup\fontsize{8.5}{10}\selectfont
\setlength{\tabcolsep}{4pt}\renewcommand{\arraystretch}{1.06}
\begin{adjustbox}{max width=\textwidth,max totalheight=0.64\textheight}
% Generated from Paper/sections/regression_local_geography_robustness_11_19_reference_table.tex
\begin{tabular}{lcccc}
\toprule
    & \multicolumn{4}{c}{Dependent variable: log real hourly labor income} \\
    \cmidrule(lr){2-5}
    & (1) & (2) & (3) & (4) \\
    \midrule
Firm size: Solo & -0.302*** \sereg{0.018} & -0.272*** \sereg{0.017} & -0.267*** \sereg{0.017} & -0.272*** \sereg{0.015} \\
Firm size: 2-3 & -0.114*** \sereg{0.013} & -0.104*** \sereg{0.011} & -0.100*** \sereg{0.011} & -0.169*** \sereg{0.009} \\
Firm size: 4-5 & 0.005 \sereg{0.008} & 0.005 \sereg{0.008} & 0.006 \sereg{0.008} & -0.047*** \sereg{0.007} \\
Firm size: 6-10 & 0.014*** \sereg{0.005} & 0.010** \sereg{0.005} & 0.011** \sereg{0.005} & -0.007 \sereg{0.004} \\
Firm size: 20-30 & -0.031*** \sereg{0.004} & -0.031*** \sereg{0.004} & -0.030*** \sereg{0.004} & -0.014*** \sereg{0.004} \\
Firm size: 31-50 & -0.031*** \sereg{0.005} & -0.027*** \sereg{0.005} & -0.027*** \sereg{0.005} & -0.008* \sereg{0.004} \\
Firm size: 51-100 & -0.014* \sereg{0.007} & -0.010 \sereg{0.007} & -0.010 \sereg{0.007} & 0.010 \sereg{0.007} \\
Firm size: 101+ & 0.057*** \sereg{0.014} & 0.066*** \sereg{0.015} & 0.058*** \sereg{0.015} & 0.066*** \sereg{0.013} \\
\midrule
    Worker and formality controls & Yes & Yes & Yes & Yes \\
Occupational-position controls & No & No & No & Yes \\
Fixed effects & $s\times t$ & $s\times t$, $d\times t$ & $s\times d\times t$ & $s\times d\times t$ \\
Observations & 5,268,514 & 5,268,514 & 5,268,375 & 5,268,375 \\
$R^2$ & 0.386 & 0.404 & 0.414 & 0.429\\
\bottomrule\end{tabular}

\end{adjustbox}\endgroup\par\vspace{0.10cm}
{\fontsize{6.5}{7.5}\selectfont\raggedright Comparison specification. Common-premium model with 11--19 as the reference. All columns include worker characteristics and formality. FE: (1) $s\times t$, (2) $s\times t$ and $d\times t$, (3)--(4) $s\times d\times t$. Column (4) adds status in employment. SEs cluster by sector-department. SEs in parentheses. * $p<0.10$, ** $p<0.05$, *** $p<0.01$.\par}
\note{Specifications as in Table~6, with 11--19 workers as the omitted category. $s$, $d$, and $t$ denote sector, department, and year. Parentheses contain sector-department-clustered standard errors. * $p<0.10$, ** $p<0.05$, *** $p<0.01$.}
\end{frame}

\begin{frame}{Paper Table 23 (1/2)}
\begin{minipage}{\textwidth}\small\textbf{Firm-size wage premiums by formality: reference category 101+}\end{minipage}\par\vspace{0.15cm}
\centering\begingroup\fontsize{8.5}{10}\selectfont
\setlength{\tabcolsep}{4pt}\renewcommand{\arraystretch}{1.06}
\begin{adjustbox}{max width=\textwidth,max totalheight=0.64\textheight}
% Generated from Paper/sections/regression_formality_interaction_coefficients_table.tex
\begin{tabular}{lccc}
\toprule
    Firm size & $\hat{\beta}_g$ & $\hat{\theta}_g$ & $\hat{\beta}_g+\hat{\theta}_g$ \\
    \midrule
2-3 & 0.197*** \sereg{0.026} & -0.175** \sereg{0.062} & 0.022 \sereg{0.067} \\
4-5 & 0.343*** \sereg{0.027} & -0.318*** \sereg{0.085} & 0.025 \sereg{0.089} \\
6-10 & 0.392*** \sereg{0.032} & -0.369*** \sereg{0.081} & 0.023 \sereg{0.086} \\
11-19 & 0.408*** \sereg{0.037} & -0.373*** \sereg{0.076} & 0.035 \sereg{0.080} \\
20-30 & 0.401*** \sereg{0.035} & -0.378*** \sereg{0.072} & 0.024 \sereg{0.079} \\
31-50 & 0.399*** \sereg{0.039} & -0.358*** \sereg{0.066} & 0.042 \sereg{0.076} \\
51-100 & 0.409*** \sereg{0.044} & -0.340*** \sereg{0.067} & 0.069 \sereg{0.075} \\
101+ & 0.336*** \sereg{0.046} & -0.182** \sereg{0.083} & 0.154 \sereg{0.088}\\
\bottomrule\end{tabular}

\end{adjustbox}\endgroup\par\vspace{0.10cm}
{\fontsize{6.5}{7.5}\selectfont\raggedright Within-formality log contrasts relative to 101+ workers. All columns include worker controls and GEIH weights. FE: (1) $s\times t$, (2) $s\times t$ and $d\times t$, (3)--(4) $s\times d\times t$. Column (4) adds status in employment and is preferred. SEs cluster by sector-department. SEs in parentheses. * $p<0.10$, ** $p<0.05$, *** $p<0.01$.\par}
\note{Within-formality log-wage contrasts relative to 101+ workers, from the specifications in Table~7. $s$, $d$, and $t$ denote sector, department, and year. Parentheses contain sector-department-clustered standard errors. * $p<0.10$, ** $p<0.05$, *** $p<0.01$.}
\end{frame}

\begin{frame}{Paper Table 23 (2/2)}
\begin{minipage}{\textwidth}\small\textbf{Firm-size wage premiums by formality: reference category 101+}\end{minipage}\par\vspace{0.15cm}
\centering\begingroup\fontsize{8.5}{10}\selectfont
\setlength{\tabcolsep}{4pt}\renewcommand{\arraystretch}{1.06}
\begin{adjustbox}{max width=\textwidth,max totalheight=0.64\textheight}
% Generated from Paper/sections/regression_local_geography_formality_robustness_101_reference_table.tex
\begin{tabular}{lcccc}
\toprule
    & \multicolumn{4}{c}{Dependent variable: log real hourly labor income} \\
    \cmidrule(lr){2-5}
    & (1) & (2) & (3) & \textbf{(4)} \\
    \midrule
\multicolumn{5}{l}{\textit{Formal workers}} \\
Firm size: Solo & -0.167*** \sereg{0.038} & -0.174*** \sereg{0.039} & -0.160*** \sereg{0.039} & -0.159*** \sereg{0.037} \\
Firm size: 2-3 & -0.145*** \sereg{0.018} & -0.149*** \sereg{0.018} & -0.138*** \sereg{0.018} & -0.234*** \sereg{0.018} \\
Firm size: 4-5 & -0.137*** \sereg{0.012} & -0.142*** \sereg{0.013} & -0.132*** \sereg{0.012} & -0.182*** \sereg{0.012} \\
Firm size: 6-10 & -0.138*** \sereg{0.012} & -0.144*** \sereg{0.013} & -0.134*** \sereg{0.012} & -0.146*** \sereg{0.011} \\
Firm size: 11-19 & -0.124*** \sereg{0.015} & -0.129*** \sereg{0.016} & -0.119*** \sereg{0.016} & -0.118*** \sereg{0.014} \\
Firm size: 20-30 & -0.133*** \sereg{0.014} & -0.139*** \sereg{0.015} & -0.129*** \sereg{0.015} & -0.116*** \sereg{0.012} \\
Firm size: 31-50 & -0.114*** \sereg{0.013} & -0.118*** \sereg{0.014} & -0.109*** \sereg{0.014} & -0.095*** \sereg{0.012} \\
Firm size: 51-100 & -0.088*** \sereg{0.012} & -0.093*** \sereg{0.013} & -0.084*** \sereg{0.013} & -0.070*** \sereg{0.011} \\
\midrule
    Worker controls & Yes & Yes & Yes & Yes \\
Status-in-employment controls & No & No & No & Yes \\
Fixed effects & $s\times t$ & $s\times t$, $d\times t$ & $s\times d\times t$ & $s\times d\times t$ \\
Observations & 5,426,316 & 5,426,316 & 5,426,174 & 5,426,174 \\
$R^2$ & 0.394 & 0.412 & 0.422 & 0.437\\
\bottomrule\end{tabular}

\end{adjustbox}\endgroup\par\vspace{0.10cm}
{\fontsize{6.5}{7.5}\selectfont\raggedright Within-formality log contrasts relative to 101+ workers. All columns include worker controls and GEIH weights. FE: (1) $s\times t$, (2) $s\times t$ and $d\times t$, (3)--(4) $s\times d\times t$. Column (4) adds status in employment and is preferred. SEs cluster by sector-department. SEs in parentheses. * $p<0.10$, ** $p<0.05$, *** $p<0.01$.\par}
\note{Within-formality log-wage contrasts relative to 101+ workers, from the specifications in Table~7. $s$, $d$, and $t$ denote sector, department, and year. Parentheses contain sector-department-clustered standard errors. * $p<0.10$, ** $p<0.05$, *** $p<0.01$.}
\end{frame}

\begin{frame}{Paper Table 24 (1/2)}
\begin{minipage}{\textwidth}\small\textbf{Firm-size wage premiums by formality: reference category 11--19}\end{minipage}\par\vspace{0.15cm}
\centering\begingroup\fontsize{8.5}{10}\selectfont
\setlength{\tabcolsep}{4pt}\renewcommand{\arraystretch}{1.06}
\begin{adjustbox}{max width=\textwidth,max totalheight=0.64\textheight}
% Generated from Paper/sections/regression_local_geography_formality_robustness_11_19_reference_table.tex
\begin{tabular}{lcccc}
\toprule
    & \multicolumn{4}{c}{Dependent variable: log real hourly labor income} \\
    \cmidrule(lr){2-5}
    & (1) & (2) & (3) & \textbf{(4)} \\
    \midrule
\multicolumn{5}{l}{\textit{Informal workers}} \\
Firm size: Solo & -0.415*** \sereg{0.016} & -0.375*** \sereg{0.014} & -0.369*** \sereg{0.014} & -0.354*** \sereg{0.018} \\
Firm size: 2-3 & -0.212*** \sereg{0.014} & -0.192*** \sereg{0.011} & -0.187*** \sereg{0.011} & -0.242*** \sereg{0.011} \\
Firm size: 4-5 & -0.063*** \sereg{0.009} & -0.055*** \sereg{0.008} & -0.052*** \sereg{0.009} & -0.098*** \sereg{0.007} \\
Firm size: 6-10 & -0.016*** \sereg{0.006} & -0.016*** \sereg{0.006} & -0.015** \sereg{0.006} & -0.030*** \sereg{0.006} \\
Firm size: 20-30 & -0.005 \sereg{0.007} & -0.007 \sereg{0.007} & -0.007 \sereg{0.007} & 0.010 \sereg{0.007} \\
Firm size: 31-50 & -0.009 \sereg{0.011} & -0.009 \sereg{0.011} & -0.012 \sereg{0.011} & 0.009 \sereg{0.010} \\
Firm size: 51-100 & 0.004 \sereg{0.014} & 0.002 \sereg{0.013} & -0.002 \sereg{0.013} & 0.021 \sereg{0.013} \\
Firm size: 101+ & -0.074*** \sereg{0.016} & -0.066*** \sereg{0.016} & -0.077*** \sereg{0.016} & -0.040** \sereg{0.017} \\
\midrule
    Worker controls & Yes & Yes & Yes & Yes \\
Status-in-employment controls & No & No & No & Yes \\
Fixed effects & $s\times t$ & $s\times t$, $d\times t$ & $s\times d\times t$ & $s\times d\times t$ \\
Observations & 5,426,316 & 5,426,316 & 5,426,174 & 5,426,174 \\
$R^2$ & 0.394 & 0.412 & 0.422 & 0.437\\
\bottomrule\end{tabular}

\end{adjustbox}\endgroup\par\vspace{0.10cm}
{\fontsize{6.5}{7.5}\selectfont\raggedright Within-formality log contrasts relative to 11--19 workers. All columns include worker controls and GEIH weights. FE: (1) $s\times t$, (2) $s\times t$ and $d\times t$, (3)--(4) $s\times d\times t$. Column (4) adds status in employment and is preferred. SEs cluster by sector-department. SEs in parentheses. * $p<0.10$, ** $p<0.05$, *** $p<0.01$.\par}
\note{Within-formality log-wage contrasts relative to 11--19 workers, from the specifications in Table~7. $s$, $d$, and $t$ denote sector, department, and year. Parentheses contain sector-department-clustered standard errors. * $p<0.10$, ** $p<0.05$, *** $p<0.01$.}
\end{frame}

\begin{frame}{Paper Table 24 (2/2)}
\begin{minipage}{\textwidth}\small\textbf{Firm-size wage premiums by formality: reference category 11--19}\end{minipage}\par\vspace{0.15cm}
\centering\begingroup\fontsize{8.5}{10}\selectfont
\setlength{\tabcolsep}{4pt}\renewcommand{\arraystretch}{1.06}
\begin{adjustbox}{max width=\textwidth,max totalheight=0.64\textheight}
% Generated from Paper/sections/regression_local_geography_formality_robustness_11_19_reference_table.tex
\begin{tabular}{lcccc}
\toprule
    & \multicolumn{4}{c}{Dependent variable: log real hourly labor income} \\
    \cmidrule(lr){2-5}
    & (1) & (2) & (3) & \textbf{(4)} \\
    \midrule
\multicolumn{5}{l}{\textit{Formal workers}} \\
Firm size: Solo & -0.043 \sereg{0.034} & -0.045 \sereg{0.034} & -0.041 \sereg{0.034} & -0.041 \sereg{0.031} \\
Firm size: 2-3 & -0.021 \sereg{0.014} & -0.021 \sereg{0.013} & -0.019 \sereg{0.013} & -0.116*** \sereg{0.014} \\
Firm size: 4-5 & -0.013 \sereg{0.009} & -0.013 \sereg{0.008} & -0.013 \sereg{0.008} & -0.064*** \sereg{0.007} \\
Firm size: 6-10 & -0.014* \sereg{0.007} & -0.016** \sereg{0.007} & -0.015** \sereg{0.007} & -0.028*** \sereg{0.007} \\
Firm size: 20-30 & -0.009** \sereg{0.004} & -0.010** \sereg{0.005} & -0.010** \sereg{0.004} & 0.002 \sereg{0.004} \\
Firm size: 31-50 & 0.010* \sereg{0.005} & 0.011** \sereg{0.005} & 0.010** \sereg{0.005} & 0.023*** \sereg{0.005} \\
Firm size: 51-100 & 0.036*** \sereg{0.007} & 0.035*** \sereg{0.007} & 0.035*** \sereg{0.007} & 0.048*** \sereg{0.006} \\
Firm size: 101+ & 0.124*** \sereg{0.015} & 0.129*** \sereg{0.016} & 0.119*** \sereg{0.016} & 0.118*** \sereg{0.014} \\
\midrule
    Worker controls & Yes & Yes & Yes & Yes \\
Status-in-employment controls & No & No & No & Yes \\
Fixed effects & $s\times t$ & $s\times t$, $d\times t$ & $s\times d\times t$ & $s\times d\times t$ \\
Observations & 5,426,316 & 5,426,316 & 5,426,174 & 5,426,174 \\
$R^2$ & 0.394 & 0.412 & 0.422 & 0.437\\
\bottomrule\end{tabular}

\end{adjustbox}\endgroup\par\vspace{0.10cm}
{\fontsize{6.5}{7.5}\selectfont\raggedright Within-formality log contrasts relative to 11--19 workers. All columns include worker controls and GEIH weights. FE: (1) $s\times t$, (2) $s\times t$ and $d\times t$, (3)--(4) $s\times d\times t$. Column (4) adds status in employment and is preferred. SEs cluster by sector-department. SEs in parentheses. * $p<0.10$, ** $p<0.05$, *** $p<0.01$.\par}
\note{Within-formality log-wage contrasts relative to 11--19 workers, from the specifications in Table~7. $s$, $d$, and $t$ denote sector, department, and year. Parentheses contain sector-department-clustered standard errors. * $p<0.10$, ** $p<0.05$, *** $p<0.01$.}
\end{frame}

\begin{frame}{Paper Table 25}
\begin{minipage}{\textwidth}\small\textbf{Formality-specific coefficients in the sector-year comparison model}\end{minipage}\par\vspace{0.15cm}
\centering\begingroup\fontsize{8.5}{10}\selectfont
\setlength{\tabcolsep}{4pt}\renewcommand{\arraystretch}{1.06}
\begin{adjustbox}{max width=\textwidth,max totalheight=0.64\textheight}
% Generated from Paper/sections/regression_formality_interaction_coefficients_table.tex
\begin{tabular}{lccc}
\toprule
    Firm size & $\hat{\beta}_g$ & $\hat{\theta}_g$ & $\hat{\beta}_g+\hat{\theta}_g$ \\
    \midrule
2-3 & 0.203*** \sereg{0.028} & -0.181** \sereg{0.062} & 0.022 \sereg{0.067} \\
4-5 & 0.352*** \sereg{0.029} & -0.322*** \sereg{0.084} & 0.030 \sereg{0.089} \\
6-10 & 0.399*** \sereg{0.032} & -0.369*** \sereg{0.081} & 0.029 \sereg{0.087} \\
11-19 & 0.415*** \sereg{0.038} & -0.372*** \sereg{0.075} & 0.043 \sereg{0.079} \\
20-30 & 0.410*** \sereg{0.037} & -0.376*** \sereg{0.071} & 0.034 \sereg{0.079} \\
31-50 & 0.406*** \sereg{0.041} & -0.353*** \sereg{0.067} & 0.053 \sereg{0.076} \\
51-100 & 0.419*** \sereg{0.047} & -0.339*** \sereg{0.067} & 0.079 \sereg{0.074} \\
101+ & 0.341*** \sereg{0.047} & -0.174* \sereg{0.082} & 0.167* \sereg{0.087}\\
\bottomrule\end{tabular}

\end{adjustbox}\endgroup\par\vspace{0.10cm}
{\fontsize{6.5}{7.5}\selectfont\raggedright Comparison specification. Log-point coefficients from Equation~(9), with solo workers as the reference within each formality group and tildes suppressed. The first and third columns are informal and formal size gaps. The middle column is the formal-minus-informal difference in those size gaps. Parentheses contain sector-clustered standard errors. * $p<0.10$, ** $p<0.05$, *** $p<0.01$.\par}
\note{Log-point coefficients from Equation~(9), with solo workers as the reference within each formality group and tildes suppressed. The first and third columns are informal and formal size gaps. The middle column is the formal-minus-informal difference in those size gaps. Parentheses contain sector-clustered standard errors. * $p<0.10$, ** $p<0.05$, *** $p<0.01$.}
\end{frame}

\begin{frame}{Paper Table 26}
\begin{minipage}{\textwidth}\small\textbf{Formal 101+ premiums in the sector-year comparison model}\end{minipage}\par\vspace{0.15cm}
\centering\begingroup\fontsize{8.5}{10}\selectfont
\setlength{\tabcolsep}{4pt}\renewcommand{\arraystretch}{1.06}
\begin{adjustbox}{max width=\textwidth,max totalheight=0.64\textheight}
% Generated from Paper/sections/regression_formality_interaction_coefficients_101_reference.tex
\begin{tabular}{lccc}
\toprule
    Firm size & $\hat{\beta}_g$ & $\hat{\theta}_g$ & $\hat{\beta}_g+\hat{\theta}_g$ \\
    \midrule
Solo & -0.336*** \sereg{0.046} & 0.182** \sereg{0.083} & -0.154 \sereg{0.088} \\
2-3 & -0.139*** \sereg{0.034} & 0.007 \sereg{0.047} & -0.132*** \sereg{0.040} \\
4-5 & 0.007 \sereg{0.031} & -0.136** \sereg{0.048} & -0.129*** \sereg{0.033} \\
6-10 & 0.056* \sereg{0.029} & -0.187*** \sereg{0.046} & -0.131*** \sereg{0.035} \\
11-19 & 0.072** \sereg{0.027} & -0.191*** \sereg{0.044} & -0.119** \sereg{0.041} \\
20-30 & 0.066** \sereg{0.024} & -0.196*** \sereg{0.042} & -0.130*** \sereg{0.040} \\
31-50 & 0.064** \sereg{0.022} & -0.175*** \sereg{0.041} & -0.112*** \sereg{0.036} \\
51-100 & 0.073*** \sereg{0.021} & -0.158*** \sereg{0.039} & -0.085** \sereg{0.030}\\
\bottomrule\end{tabular}

\end{adjustbox}\endgroup\par\vspace{0.10cm}
{\fontsize{6.5}{7.5}\selectfont\raggedright Comparison specification. Each premium compares formal 101+ workers with formal workers in the listed reference category, using Equation~(9). Two-sided $p$-values use sector-clustered standard errors and 16 degrees of freedom.\par}
\note{Each premium compares formal 101+ workers with formal workers in the listed reference category, using Equation~(9). Two-sided $p$-values use sector-clustered standard errors and 16 degrees of freedom.}
\end{frame}

\begin{frame}{Paper Table 27}
\begin{minipage}{\textwidth}\small\textbf{Formal-informal wage gaps in the sector-year comparison model}\end{minipage}\par\vspace{0.15cm}
\centering\begingroup\fontsize{8.5}{10}\selectfont
\setlength{\tabcolsep}{4pt}\renewcommand{\arraystretch}{1.06}
\begin{adjustbox}{max width=\textwidth,max totalheight=0.64\textheight}
% Generated from Paper/sections/regression_formality_gap_by_size_table.tex
\begin{tabular}{lccc}
\toprule
    Firm size & Log gap & Premium & $p$-value \\
    \midrule
Solo & 0.526*** \sereg{0.062} & 69.2\% & $<0.001$ \\
2-3 & 0.345*** \sereg{0.017} & 41.2\% & $<0.001$ \\
4-5 & 0.204*** \sereg{0.035} & 22.6\% & $<0.001$ \\
6-10 & 0.157*** \sereg{0.031} & 16.9\% & $<0.001$ \\
11-19 & 0.154*** \sereg{0.024} & 16.7\% & $<0.001$ \\
20-30 & 0.150*** \sereg{0.018} & 16.2\% & $<0.001$ \\
31-50 & 0.173*** \sereg{0.015} & 18.8\% & $<0.001$ \\
51-100 & 0.186*** \sereg{0.019} & 20.5\% & $<0.001$ \\
101+ & 0.352*** \sereg{0.044} & 42.2\% & $<0.001$\\
\bottomrule\end{tabular}

\end{adjustbox}\endgroup\par\vspace{0.10cm}
{\fontsize{6.5}{7.5}\selectfont\raggedright Comparison specification. Formal-minus-informal contrasts from Equation~(9) at each firm size. Percentage premiums use informal workers as the denominator. Parentheses contain sector-clustered standard errors. * $p<0.10$, ** $p<0.05$, *** $p<0.01$.\par}
\note{Formal-minus-informal contrasts from Equation~(9) at each firm size. Percentage premiums use informal workers as the denominator. Parentheses contain sector-clustered standard errors. * $p<0.10$, ** $p<0.05$, *** $p<0.01$.}
\end{frame}

\begin{frame}{Paper Table 28}
\begin{minipage}{\textwidth}\small\textbf{Firm-size coefficients by formality status: 101+ workers as reference}\end{minipage}\par\vspace{0.15cm}
\centering\begingroup\fontsize{8.5}{10}\selectfont
\setlength{\tabcolsep}{4pt}\renewcommand{\arraystretch}{1.06}
\begin{adjustbox}{max width=\textwidth,max totalheight=0.64\textheight}
% Generated from Paper/sections/regression_formality_interaction_coefficients_101_reference.tex
\begin{tabular}{lccc}
\toprule
    Firm size & $\hat{\beta}_g$ & $\hat{\theta}_g$ & $\hat{\beta}_g+\hat{\theta}_g$ \\
    \midrule
Solo & -0.341*** \sereg{0.047} & 0.174* \sereg{0.082} & -0.167* \sereg{0.087} \\
2-3 & -0.138*** \sereg{0.035} & -0.007 \sereg{0.047} & -0.145*** \sereg{0.040} \\
4-5 & 0.011 \sereg{0.032} & -0.148*** \sereg{0.049} & -0.137*** \sereg{0.033} \\
6-10 & 0.058* \sereg{0.030} & -0.195*** \sereg{0.047} & -0.138*** \sereg{0.034} \\
11-19 & 0.074** \sereg{0.027} & -0.198*** \sereg{0.044} & -0.124*** \sereg{0.040} \\
20-30 & 0.069*** \sereg{0.024} & -0.202*** \sereg{0.041} & -0.133*** \sereg{0.039} \\
31-50 & 0.065** \sereg{0.023} & -0.179*** \sereg{0.041} & -0.114*** \sereg{0.035} \\
51-100 & 0.078*** \sereg{0.022} & -0.166*** \sereg{0.040} & -0.088*** \sereg{0.029}\\
\bottomrule\end{tabular}

\end{adjustbox}\endgroup\par\vspace{0.10cm}
{\fontsize{6.5}{7.5}\selectfont\raggedright Comparison specification. Log-point coefficients from Equation~(9), with 101+ workers as the reference within each formality group and tildes suppressed. The first and third columns are informal and formal size gaps. The middle column is the formal-minus-informal difference in those size gaps. Parentheses contain sector-clustered standard errors. * $p<0.10$, ** $p<0.05$, *** $p<0.01$.\par}
\note{Log-point coefficients from Equation~(9), with 101+ workers as the reference within each formality group and tildes suppressed. The first and third columns are informal and formal size gaps. The middle column is the formal-minus-informal difference in those size gaps. Parentheses contain sector-clustered standard errors. * $p<0.10$, ** $p<0.05$, *** $p<0.01$.}
\end{frame}

\begin{frame}{Paper Table 29}
\begin{minipage}{\textwidth}\small\textbf{Firm-size coefficients by formality status: 11--19 workers as reference}\end{minipage}\par\vspace{0.15cm}
\centering\begingroup\fontsize{8.5}{10}\selectfont
\setlength{\tabcolsep}{4pt}\renewcommand{\arraystretch}{1.06}
\begin{adjustbox}{max width=\textwidth,max totalheight=0.64\textheight}
% Generated from Paper/sections/regression_formality_interaction_coefficients_11_19_reference.tex
\begin{tabular}{lccc}
\toprule
    Firm size & $\hat{\beta}_g$ & $\hat{\theta}_g$ & $\hat{\beta}_g+\hat{\theta}_g$ \\
    \midrule
Solo & -0.415*** \sereg{0.038} & 0.372*** \sereg{0.075} & -0.043 \sereg{0.079} \\
2-3 & -0.212*** \sereg{0.025} & 0.191*** \sereg{0.022} & -0.021 \sereg{0.024} \\
4-5 & -0.063*** \sereg{0.019} & 0.049** \sereg{0.018} & -0.013 \sereg{0.018} \\
6-10 & -0.016 \sereg{0.010} & 0.002 \sereg{0.012} & -0.014 \sereg{0.012} \\
20-30 & -0.005 \sereg{0.007} & -0.004 \sereg{0.009} & -0.009* \sereg{0.005} \\
31-50 & -0.009 \sereg{0.009} & 0.019 \sereg{0.012} & 0.010 \sereg{0.008} \\
51-100 & 0.004 \sereg{0.023} & 0.032 \sereg{0.026} & 0.036** \sereg{0.015} \\
101+ & -0.074** \sereg{0.027} & 0.198*** \sereg{0.044} & 0.124*** \sereg{0.040}\\
\bottomrule\end{tabular}

\end{adjustbox}\endgroup\par\vspace{0.10cm}
{\fontsize{6.5}{7.5}\selectfont\raggedright Comparison specification. Log-point coefficients from Equation~(9), with 11--19 workers as the reference within each formality group and tildes suppressed. The first and third columns are informal and formal size gaps. The middle column is the formal-minus-informal difference in those size gaps. Parentheses contain sector-clustered standard errors. * $p<0.10$, ** $p<0.05$, *** $p<0.01$.\par}
\note{Log-point coefficients from Equation~(9), with 11--19 workers as the reference within each formality group and tildes suppressed. The first and third columns are informal and formal size gaps. The middle column is the formal-minus-informal difference in those size gaps. Parentheses contain sector-clustered standard errors. * $p<0.10$, ** $p<0.05$, *** $p<0.01$.}
\end{frame}

\begin{frame}{Paper Table 30}
\begin{minipage}{\textwidth}\small\textbf{Testing whether the firm-size wage premium changed over time}\end{minipage}\par\vspace{0.15cm}
\centering\begingroup\fontsize{8.5}{10}\selectfont
\setlength{\tabcolsep}{4pt}\renewcommand{\arraystretch}{1.06}
\begin{adjustbox}{max width=\textwidth,max totalheight=0.64\textheight}
% Generated from Paper/sections/regression_firm_size_trend_test_11_19_reference.tex
\begin{tabular}{lcccc}
\toprule
    Firm size & Annual trend & S.E. & $p$-value & Implied 2008--2025 change \\
    \midrule
Solo & -0.0027 & (0.0039) & 0.491 & -4.5\% \\
2-3 & 0.0005 & (0.0016) & 0.746 & 0.9\% \\
4-5 & 0.0005 & (0.0014) & 0.736 & 0.8\% \\
6-10 & -0.0006 & (0.0011) & 0.611 & -1.0\% \\
20-30 & -0.0009 & (0.0010) & 0.392 & -1.5\% \\
31-50 & 0.0000 & (0.0012) & 0.988 & 0.0\% \\
51-100 & -0.0001 & (0.0014) & 0.940 & -0.2\% \\
101+ & 0.0017 & (0.0012) & 0.175 & 3.0\%\\
\bottomrule\end{tabular}

\end{adjustbox}\endgroup\par\vspace{0.10cm}
{\fontsize{6.5}{7.5}\selectfont\raggedright Comparison specification. Annual changes in log-wage gaps from Equation~(11), expressed relative to solo workers. The final column is the compounded 2008--2025 percentage change in that wage ratio, not a percentage-point change in the premium. Two-sided tests of zero trends use sector-clustered standard errors.\par}
\note{Annual changes in log-wage gaps from Equation~(11), expressed relative to solo workers. The final column is the compounded 2008--2025 percentage change in that wage ratio, not a percentage-point change in the premium. Two-sided tests of zero trends use sector-clustered standard errors.}
\end{frame}

\begin{frame}{Paper Table 31}
\begin{minipage}{\textwidth}\small\textbf{Testing whether wage gaps relative to 101+ workers changed over time}\end{minipage}\par\vspace{0.15cm}
\centering\begingroup\fontsize{8.5}{10}\selectfont
\setlength{\tabcolsep}{4pt}\renewcommand{\arraystretch}{1.06}
\begin{adjustbox}{max width=\textwidth,max totalheight=0.64\textheight}
% Generated from Paper/sections/regression_firm_size_trend_test_101_reference.tex
\begin{tabular}{lcccc}
\toprule
    Firm size & Annual trend & S.E. & $p$-value & Implied 2008--2025 change \\
    \midrule
Solo & -0.0038 & (0.0041) & 0.372 & -6.3\% \\
2-3 & -0.0012 & (0.0024) & 0.608 & -2.1\% \\
4-5 & -0.0008 & (0.0023) & 0.733 & -1.4\% \\
6-10 & -0.0016 & (0.0017) & 0.355 & -2.7\% \\
11-19 & -0.0012 & (0.0012) & 0.337 & -2.0\% \\
20-30 & -0.0024 & (0.0009) & 0.020 & -4.0\% \\
31-50 & -0.0014 & (0.0008) & 0.114 & -2.3\% \\
51-100 & -0.0013 & (0.0013) & 0.327 & -2.1\%\\
\bottomrule\end{tabular}

\end{adjustbox}\endgroup\par\vspace{0.10cm}
{\fontsize{6.5}{7.5}\selectfont\raggedright Comparison specification. Annual changes in log-wage gaps from Equation~(11), expressed relative to 101+ workers. The final column is the compounded 2008--2025 percentage change in that wage ratio, not a percentage-point change in the premium. Two-sided tests of zero trends use sector-clustered standard errors.\par}
\note{Annual changes in log-wage gaps from Equation~(11), expressed relative to 101+ workers. The final column is the compounded 2008--2025 percentage change in that wage ratio, not a percentage-point change in the premium. Two-sided tests of zero trends use sector-clustered standard errors.}
\end{frame}

\begin{frame}{Paper Table 32}
\begin{minipage}{\textwidth}\small\textbf{Testing whether wage gaps relative to 11--19 workers changed over time}\end{minipage}\par\vspace{0.15cm}
\centering\begingroup\fontsize{8.5}{10}\selectfont
\setlength{\tabcolsep}{4pt}\renewcommand{\arraystretch}{1.06}
\begin{adjustbox}{max width=\textwidth,max totalheight=0.64\textheight}
% Generated from Paper/sections/regression_preferred_quadratic_table.tex
\begin{tabular}{lcc}
\toprule
    Term & Formal workers & Informal workers \\[0.25em]
    \midrule
Log assigned size ($\widehat{a}_{1f}$) & \shortstack{-0.0095\\\sereg{0.0148}} & \shortstack{0.1947*** \\
\sereg{0.0133}} \\
Squared log assigned size ($\widehat{a}_{2f}$) & \shortstack{0.0088***\\\sereg{0.0021}} & \shortstack{-0.0250***\\\sereg{0.0025}} \\
\midrule
    Implied 101+/solo premium (\%) & 18.8 & 41.8\\
\bottomrule\end{tabular}

\end{adjustbox}\endgroup\par\vspace{0.10cm}
{\fontsize{6.5}{7.5}\selectfont\raggedright Comparison specification. Annual changes in log-wage gaps from Equation~(11), expressed relative to 11--19 workers. The final column is the compounded 2008--2025 percentage change in that wage ratio, not a percentage-point change in the premium. Two-sided tests of zero trends use sector-clustered standard errors.\par}
\note{Annual changes in log-wage gaps from Equation~(11), expressed relative to 11--19 workers. The final column is the compounded 2008--2025 percentage change in that wage ratio, not a percentage-point change in the premium. Two-sided tests of zero trends use sector-clustered standard errors.}
\end{frame}

\begin{frame}{Paper Table 33}
\begin{minipage}{\textwidth}\small\textbf{Benchmarking Colombia's approximate firm-size wage elasticity}\end{minipage}\par\vspace{0.15cm}
\centering\begingroup\fontsize{8.5}{10}\selectfont
\setlength{\tabcolsep}{4pt}\renewcommand{\arraystretch}{1.06}
\begin{adjustbox}{max width=\textwidth,max totalheight=0.64\textheight}
% Generated from Paper/sections/regression_preferred_quadratic_sensitivity_table.tex
\begin{tabular}{rrrrrr}
\toprule
    Size for 101+ & Formal (\%) & Informal (\%) & $p_F$ & $p_I$ & $p_{Joint}$ \\[0.25em]
    \midrule
\multicolumn{6}{l}{\textit{Sector-department clustering}} \\
101 & 17.8 & 43.1 & $<0.001$ & $<0.001$ & $<0.001$ \\
150 & 18.8 & 41.8 & $<0.001$ & $<0.001$ & $<0.001$ \\
250 & 19.8 & 40.4 & $<0.001$ & $<0.001$ & $<0.001$ \\
\midrule
    \multicolumn{6}{l}{\textit{Sector clustering}} \\
101 & 17.8 & 43.1 & $<0.001$ & $<0.001$ & $<0.001$ \\
150 & 18.8 & 41.8 & $<0.001$ & $<0.001$ & $<0.001$ \\
250 & 19.8 & 40.4 & $<0.001$ & 0.001 & $<0.001$\\
\bottomrule\end{tabular}

\end{adjustbox}\endgroup\par\vspace{0.10cm}
{\fontsize{6.5}{7.5}\selectfont\raggedright Comparison specification. Literature means are from Table 1 of Diegmann et al. (2026). The controlled full-sample row uses Equation~(12). Formal and informal rows re-estimate it on separate samples. Standard errors are clustered by sector. * $p<0.10$, ** $p<0.05$, *** $p<0.01$.\par}
\note{Literature means are from Table 1 of Diegmann et al. (2026). The controlled full-sample row uses Equation~(12). Formal and informal rows re-estimate it on separate samples. Standard errors are clustered by sector. * $p<0.10$, ** $p<0.05$, *** $p<0.01$.}
\end{frame}

\begin{frame}{Paper Table 34}
\begin{minipage}{\textwidth}\small\textbf{Quadratic size profiles with preferred controls}\end{minipage}\par\vspace{0.15cm}
\centering\begingroup\fontsize{8.5}{10}\selectfont
\setlength{\tabcolsep}{4pt}\renewcommand{\arraystretch}{1.06}
\begin{adjustbox}{max width=\textwidth,max totalheight=0.64\textheight}
% Generated from Paper/sections/regression_preferred_quadratic_table.tex
\begin{tabular}{lcc}
\toprule
    Term & Formal workers & Informal workers \\[0.25em]
    \midrule
Log assigned size ($\widehat{a}_{1f}$) & \shortstack{-0.0086\\\sereg{0.0147}} & \shortstack{0.1956*** \\
\sereg{0.0134}} \\
Squared log assigned size ($\widehat{a}_{2f}$) & \shortstack{0.0089***\\\sereg{0.0021}} & \shortstack{-0.0251***\\\sereg{0.0026}} \\
\midrule
    Implied 101+/solo premium (\%) & 19.8 & 41.9\\
\bottomrule\end{tabular}

\end{adjustbox}\endgroup\par\vspace{0.10cm}
{\fontsize{6.5}{7.5}\selectfont\raggedright Both columns come from Equation~(13) on the preferred sample, with 101+ assigned 150 workers. The last row reports percentage premiums relative to solo work within each formality group. Parentheses contain sector-department-clustered standard errors. * $p<0.10$, ** $p<0.05$, *** $p<0.01$.\par}
\note{Both columns come from Equation~(13) on the preferred sample, with 101+ assigned 150 workers. The last row reports percentage premiums relative to solo work within each formality group. Parentheses contain sector-department-clustered standard errors. * $p<0.10$, ** $p<0.05$, *** $p<0.01$.}
\end{frame}

\begin{frame}{Paper Table 35}
\begin{minipage}{\textwidth}\small\textbf{Top-bin sensitivity and tests of the quadratic restriction}\end{minipage}\par\vspace{0.15cm}
\centering\begingroup\fontsize{8.5}{10}\selectfont
\setlength{\tabcolsep}{4pt}\renewcommand{\arraystretch}{1.06}
\begin{adjustbox}{max width=\textwidth,max totalheight=0.64\textheight}
% Generated from Paper/sections/regression_preferred_quadratic_sensitivity_table.tex
\begin{tabular}{rrrrrr}
\toprule
    Size for 101+ & Formal (\%) & Informal (\%) & $p_F$ & $p_I$ & $p_{Joint}$ \\[0.25em]
    \midrule
\multicolumn{6}{l}{\textit{Sector-department clustering}} \\
101 & 18.7 & 43.2 & $<0.001$ & $<0.001$ & $<0.001$ \\
150 & 19.8 & 41.9 & $<0.001$ & $<0.001$ & $<0.001$ \\
250 & 20.8 & 40.5 & $<0.001$ & $<0.001$ & $<0.001$ \\
\midrule
    \multicolumn{6}{l}{\textit{Sector clustering}} \\
101 & 18.7 & 43.2 & $<0.001$ & $<0.001$ & $<0.001$ \\
150 & 19.8 & 41.9 & $<0.001$ & $<0.001$ & $<0.001$ \\
250 & 20.8 & 40.5 & $<0.001$ & $<0.001$ & $<0.001$\\
\bottomrule\end{tabular}

\end{adjustbox}\endgroup\par\vspace{0.10cm}
{\fontsize{6.5}{7.5}\selectfont\raggedright Formal and informal percentages are the implied 101+/solo premiums from Equation~(13), re-estimated for each assigned top-bin size. All other size assignments, observations, controls, and weights remain fixed. The $p$-values test whether the unrestricted categorical profiles from Equation~(4) lie on the specified quadratic. Wald tests use an $F_{q,G-1}$ reference distribution, with $q=6$ within each formality group and $q=12$ jointly. There are 559 sector-department clusters or 17 sector clusters. The value 150 is the baseline assignment, not an observed mean for 101+.\par}
\note{Formal and informal percentages are the implied 101+/solo premiums from Equation~(13), re-estimated for each assigned top-bin size. All other size assignments, observations, controls, and weights remain fixed. The $p$-values test whether the unrestricted categorical profiles from Equation~(4) lie on the specified quadratic. Wald tests use an $F_{q,G-1}$ reference distribution, with $q=6$ within each formality group and $q=12$ jointly. There are 559 sector-department clusters or 17 sector clusters. The value 150 is the baseline assignment, not an observed mean for 101+.}
\end{frame}

\begin{frame}{Paper Table 36}
\begin{minipage}{\textwidth}\small\textbf{Income scenarios with alternative baselines and destinations}\end{minipage}\par\vspace{0.15cm}
\centering\begingroup\fontsize{8.5}{10}\selectfont
\setlength{\tabcolsep}{4pt}\renewcommand{\arraystretch}{1.06}
\begin{adjustbox}{max width=\textwidth,max totalheight=0.64\textheight}
% Generated from Paper/sections/income_counterfactual_sensitivity_table.tex
\begin{tabular}{>{\raggedright\arraybackslash}p{0.52\textwidth}cc}
\toprule
    Scenario & Predicted baseline & Observed baseline \\
    & (\%) & (\%) \\
    \midrule
Informal to formal, same size & 13.5 & 13.3 \\
Formal workers to formal 101+ & 3.8 & 3.8 \\
Informal workers to informal 101+ & 7.6 & 7.5 \\
All workers to formal 101+ & 25.9 & 25.5 \\
\midrule
    \multicolumn{3}{l}{\textit{Destinations selected from the estimated profiles}} \\
Informal workers to informal 51--100 & 10.1 & 9.9 \\
Formal to 101+, informal to 51--100 & 13.8 & 13.7\\
\bottomrule\end{tabular}

\end{adjustbox}\endgroup\par\vspace{0.10cm}
{\fontsize{6.5}{7.5}\selectfont\raggedright Percentage changes in total weekly labor income for the entire 2025 sample. The predicted baseline uses Equation~(8). The observed baseline rescales observed income by the same estimated wage ratios, preserving log residuals. Both use pooled coefficients from Equation~(4). In the last two rows, destinations maximize the estimated size component within each formality group, including moves from informal 101+ to informal 51--100.\par}
\note{Percentage changes in total weekly labor income for the entire 2025 sample. The predicted baseline uses Equation~(8). The observed baseline rescales observed income by the same estimated wage ratios, preserving log residuals. Both use pooled coefficients from Equation~(4). In the last two rows, destinations maximize the estimated size component within each formality group, including moves from informal 101+ to informal 51--100.}
\end{frame}

\begin{frame}{Paper Table 37 (1/2)}
\begin{minipage}{\textwidth}\small\textbf{Observed destination groups and sensitivity to extrapolation, 2025}\end{minipage}\par\vspace{0.15cm}
\centering\begingroup\fontsize{8.5}{10}\selectfont
\setlength{\tabcolsep}{4pt}\renewcommand{\arraystretch}{1.06}
\begin{adjustbox}{max width=\textwidth,max totalheight=0.64\textheight}
% Generated from Paper/sections/income_counterfactual_support_table.tex
\begin{tabular}{>{\raggedright\arraybackslash}p{0.40\textwidth}cccc}
\toprule
    \multicolumn{5}{l}{\textit{A. Intended moves and absent destination cells (\%)}} \\
    \midrule
Scenario & \shortstack{Intended\\movers} & \multicolumn{2}{c}{\shortstack{Absent destination\\in 2025}} & \shortstack{Absent\\in pooled} \\
\midrule
    Informal to formal, same size & 55.1 & \multicolumn{2}{c}{6.0} & 0.4 \\
Formal to 101+ & 20.2 & \multicolumn{2}{c}{6.9} & 1.8 \\
Informal to 101+ & 54.6 & \multicolumn{2}{c}{39.0} & 6.9 \\
All to formal 101+ & 75.3 & \multicolumn{2}{c}{17.2} & 4.6 \\
Informal to 51--100 & 54.9 & \multicolumn{2}{c}{55.0} & 10.8 \\
Formal to 101+, informal to 51--100 & 75.1 & \multicolumn{2}{c}{42.0} & 8.4\\
\bottomrule\end{tabular}

\end{adjustbox}\endgroup\par\vspace{0.10cm}
{\fontsize{6.5}{7.5}\selectfont\raggedright Destination cells are defined by sector, department, status in employment, size, and formality. Panel A reports intended movers as a percentage of all workers. The absent-destination columns instead use intended movers as their denominator. Pooled counts use all estimation years. Panel B permits a move only when its destination meets the stated unweighted observation count. Otherwise the worker stays at origin. The all-worker income denominator is unchanged. \textit{None} means that no intended move qualifies. For the joint selected-destination scenario, the 2025 $n\geq30$ column contains only formal moves.\par}
\note{Destination cells are defined by sector, department, status in employment, size, and formality. Panel A reports intended movers as a percentage of all workers. The absent-destination columns instead use intended movers as their denominator. Pooled counts use all estimation years. Panel B permits a move only when its destination meets the stated unweighted observation count. Otherwise the worker stays at origin. The all-worker income denominator is unchanged. \textit{None} means that no intended move qualifies. For the joint selected-destination scenario, the 2025 $n\geq30$ column contains only formal moves.}
\end{frame}

\begin{frame}{Paper Table 37 (2/2)}
\begin{minipage}{\textwidth}\small\textbf{Observed destination groups and sensitivity to extrapolation, 2025}\end{minipage}\par\vspace{0.15cm}
\centering\begingroup\fontsize{8.5}{10}\selectfont
\setlength{\tabcolsep}{4pt}\renewcommand{\arraystretch}{1.06}
\begin{adjustbox}{max width=\textwidth,max totalheight=0.64\textheight}
% Generated from Paper/sections/income_counterfactual_support_table.tex
\begin{tabular}{>{\raggedright\arraybackslash}p{0.40\textwidth}cccc}
\toprule\multicolumn{5}{l}{\textit{B. Change in total predicted income with destination screens (\%)}} \\
    \midrule
    Scenario & No screen & \shortstack{Pooled\\$n\geq30$} & \shortstack{2025\\$n\geq10$} & \shortstack{2025\\$n\geq30$} \\
    \midrule
Informal to formal, same size & 13.5 & 12.0 & 8.7 & 5.2 \\
Formal to 101+ & 3.8 & 3.2 & 2.9 & 2.4 \\
Informal to 101+ & 7.6 & 3.8 & 0.5 & \textit{None} \\
All to formal 101+ & 25.9 & 17.7 & 11.4 & 6.7 \\
Informal to 51--100 & 10.1 & 2.3 & $<0.1$ & \textit{None} \\
Formal to 101+, informal to 51--100 & 13.8 & 5.4 & 2.9 & 2.4\\
\bottomrule\end{tabular}

\end{adjustbox}\endgroup\par\vspace{0.10cm}
{\fontsize{6.5}{7.5}\selectfont\raggedright Destination cells are defined by sector, department, status in employment, size, and formality. Panel A reports intended movers as a percentage of all workers. The absent-destination columns instead use intended movers as their denominator. Pooled counts use all estimation years. Panel B permits a move only when its destination meets the stated unweighted observation count. Otherwise the worker stays at origin. The all-worker income denominator is unchanged. \textit{None} means that no intended move qualifies. For the joint selected-destination scenario, the 2025 $n\geq30$ column contains only formal moves.\par}
\note{Destination cells are defined by sector, department, status in employment, size, and formality. Panel A reports intended movers as a percentage of all workers. The absent-destination columns instead use intended movers as their denominator. Pooled counts use all estimation years. Panel B permits a move only when its destination meets the stated unweighted observation count. Otherwise the worker stays at origin. The all-worker income denominator is unchanged. \textit{None} means that no intended move qualifies. For the joint selected-destination scenario, the 2025 $n\geq30$ column contains only formal moves.}
\end{frame}
